% Gerado por tools/tei2latex.py a partir de entrada.xml. Regenerar, não editar.

%%>>> parte-1

%%>>> cap-1

%%>>> art-8

%%>>> t49

UNIÃO DA ALMA COM O CORPO

%%>>> art-9

\section*{Artigo I.}
\addcontentsline{toc}{section}{Artigo I.}

\pg{60} \textbf{União da alma com} o \textbf{corpo}

e \textbf{suas consequencias}.

43. A \textbf{alma} e o corp» constituem, pela sua união, uma só \textbf{pessoa}. --- \emph{Pessoa} é uma substancia singular, completa, racional e subsistente, de modo que é o sujeito a que se attribuem \emph{todas} as

9\emph{a} operações. A alma e o corpo, pois, constituirão uma só pes- 80a, se constituirem uma substancia singular, completa, racional e subsistente, a que se attribuam todas as operações humanas. Mas a alma e o corpo, constituindo o homem, constituem uma substancia singular, completa, racional e subsistente, a que se attribuem \emph{todas} as operações humanas, --- as intellectuaes, as sensitivas e as vegetativas. Logo a alma e o corpo constituem, pela sua união, uma só pessoa \footnote{A \emph{pessoa humana} não é só o corpo, nem só a alma, mas é o composto da alma e do corpo. O eu experimenta em si duas entidades, --- uma interior, simples e activa, --- autra exterior, extensa e passiva, e tem consciencia de que estas duas entidades formam uma só substancia completa, uma só pessoa. Isto á innegavel. Cada um de nós diz: eu vegeto, eu sinto, eu entendo, eu quero, isto é, attribue a \emph{um} e \emph{mesmo} sujeito \emph{todas} as operações humanas, e por isso reconhece a \emph{unidade} de \emph{pessoa}.\par Objectam : A \emph{pessoa} humana é só constituida pela alma. Na verdade --- 1." a alma muitas vezes diz-se \emph{homem}, e \emph{homem} denota a \emph{pessoa humana}; --- 2.º) a alma humana subsiste fóra do corpo, e por isso é \emph{supposto} ou \emph{pessoa}; --- 9.9) o corpo diz-se \emph{instrumento} da alma, e o \emph{instrumento} não constitue com o agente principal um só supposto.\par Respondemos que estes tres argumentos são insuflicientes para provarem a these dos adversarios. --- E' insufficiente o \emph{primeiro}; porque, se a alma algumas vezes se denomina \emph{homem}, esta denominação não se faz em sentido \emph{proprio}, mas em sentido \emph{tropologico}, pois toma-se a \emph{parte} pelo \emph{todo}. --- E' insufficiente o \emph{segundo}; porque a alma não subsiste \emph{fóra} do corpo como pessoa humana, mas como \emph{parte} da pessoa humana. --- E' insufficiente o \emph{terceiro} ; porque o \emph{instrumento}, se não constitue com o agente principal um\pgn{61}só supposto, quando elle é \emph{extrinseco}, constitue-o, quando é \emph{intrinseco} ; ora o corpo é instrumento \emph{intrinseco} da alma.\par Examinemos as opiniões de \emph{Descartes}, \emph{Kant}, \emph{e Rosmini}.\par a) \emph{Descartes} faz consistir o \emph{eu}, ou a \emph{pessoa humana}, exclusivamente na \emph{alma}, emquanto é pensante e distincta do corpo.\par Esta opinião é falsa. Ninguem diz: \emph{eu sou alma}; mas sim: \emph{eu sou homem}. Ora o \emph{homem} não é só a alma, mas é o composto da alma e do corpo.\par b) \emph{Kant} ensina que o \emph{eu} não é a alma, mas é a consciencia que a alma tem de si mesma, --- ou, por outras palavras: o \emph{eu} não é a substancia pensante, mas é o proprio pensamento emquanto apprehende, pela reflexão, um outro pensamento. E, como este philosopho admitte em nós um duplice ex --- 0 \emph{eu puro} e o \emph{eu empirico}, assim faz consistir o \emph{eu puro} na consciencia das formas \emph{a-priori} e o \emph{eu empirico} na consciencia das mesmas, emquanto \&pplieadas aos phenomenos sensiveis e experimentaes.\par . Tambem esta opinião é falsa. \emph{Na} verdade, o \emph{eu} é o sujeito a que se attribuem todas as propriedades. Ora a consciencia é uma faculdade e por isso, longe de ser o sujeito das propriedades, precisa de um sujeito a que se attribua e a que adhira. Logo o \emph{eu} não é a consciencia, mas é o sujeito a que a consciencia pertence. --- Accrescentamos que a distineção entre o eu \emph{puro} e o \emph{eu empirico} é gratuita e falsa. O \emph{eu} é um só, quer contemple as razões ideaes, quer perceba os phenomenos sensiveis, como tambem uma só é a consciencia, que em ambos os casos o eu tem de si mesmo. Não sendo assim, não poderiamos applicar os conceitos puros da intelligencia aos dados experimentaes, nem poderiamos ter consciencia de que em nós o sujeito que entende e que sente é um e o mesmo.\par c) \emph{Rosmini}, seguindo um caminho medio entre \emph{Descartes} e \emph{Kant}, diz que o eu não é só a alma nem só a conscencia, mas é \emph{a alma emquanto tem consciencia de si mesma}. Eis as suas palavras: « O eu é um principio activo de uma certa natureza, emquanto tem consciencia de si mesmo e pronuncia o seu acto» (\emph{Psych}., \emph{t}. 1, dei. xir). « Quando o sujeito humano, mediante as diversas operações interiores das suas faculdades, chega a adquirir a consciencia de si proprio, então \emph{torna-se eu} » (\emph{Anth}. lib. 1v. c. 4).\par Tambem esta opinião não póde acceitar-se. \emph{Rosmini}, como \emph{Kant}, confunde o eu com a percepção do \emph{eu}. O sujeito humano, pela reflexão, \emph{conhece} o \emph{eu}, mas não o \emph{forma}; e aqui não se investiga \emph{como se conhece} o \emph{eu}, mas em \emph{que consiste} o eu. --- De mais, se o \emph{eu} consistisse na alma \emph{emquanto tem consciencia} de si mesma, então não seria \emph{pessoa} --- nem a criança, que ainda não chegou á edade da reflexão, --- nem o adulto, quando actualmente não possue a consciencia de si mesmo. Ora estas consequencias, que derivam legitimamente do principio de \emph{Rosmini}, são falsas.}.

44. A \textbf{alma} e o corpo constituem, pela sua união, uma sô \textbf{natureza}. --- \emph{Natureza} é a propria essencia, emquanto é principio de operações. Logo a alma e o corpo constituirão uma só \emph{natureza}, se constituirem um só principio \emph{adequado} de operações. Mas essas duas substancias constituem, pela sua união, um só principio \emph{adequado} de operações, como o mostram as funcções da vida vegetativa e sensitiva; porquanto, o sujeito \pg{61} \pg{62}  que vegeta ou sente não é só a alma nem só o corpo, mas é o composto da alma e do corpo. Logo a alma e o corpo constituem, pela sua união, \emph{uma} só natureza \footnote{Alguns auctores, vendo que as nossas operações \emph{intellectuaes} são sempre acompanhadas pelas operações \emph{sensitivas} e produzem alterações e movimentos no proprio organismo, concluem que a alma e o corpo, mesmo em relação aos actos \emph{intellectuaes}, constituem, na sua união, \emph{uma só natureza}.\par Admittimos e sustentamos que existe uma correspondencia entre os phenomenos \emph{intellectuaes} e os \emph{organicos}; pois está averiguado que, quando pensamos, a corrente electrica, existente nos nervos, diminue de intensidade, e que uma serie de operações e transformações chimicas se produz no cerebro, nos nervos e em todo o corpo. Mas essa correspondencia, se prova que a nossa alma, nas operações \emph{intellectuaes}, precisa do concurso das faculdades \emph{sensitivas} e por isso depende \emph{extrinsecamente} do corpo, não prova nem póde provar que depende d'elle \emph{intrinsecamente}, de modo que o sujeito pensante seja o \emph{composto} da alma e de corpo; pois, como vimos, não é o \emph{composto} que pensa, mas é só a \emph{alma}. Se não é o \emph{composto} que pensa, não se póde \emph{rigorosamente} dizer que a alma e a corpo constituem, tambem em relação ás operações \emph{intellectuaes}, \emph{uma só natureza}.\par Dissemos --- \emph{rigorosamente} ; porque, se considerarmos a operação \emph{intellectual}, não como um phenomeno \emph{simples}, mas como um phenomeno \emph{complexo}, isto é, composto da operação \emph{intellectual} e das operações \emph{organicas}, que a precedem, acompanham e seguem, podemos dizer que a alma e o corpo, mesmo em relação ás operações \emph{intellectuaes}, constituem uma só natureza.\par Objectam: Se a alma e o corpo, na sua união, constituissem uma só \emph{natureza}, confundir-se-hiam. Ora esta confusão é absurda.\par Respondemos que a alma e o corpo, na sua união, constituem uma só natureza, e todavia não se confundem. Porquanto, a alma, embora communique ao corpo a subsistencia e a especie, todavia não se transforma no corpo, nem este naquella; ora seria necessaria similhante transformação, para se poder concluir que a alma e o corpo, na sua união, se confundem.\par D'onde se vé que, entre a união da nossa alma com o corpo e a união da Pessoa adoravel do Firmo DE \textsc{DEus} com a natureza humana, existe uma certa analogia e uma grande differenca. --- Existe uma certa analogia; emquanto a união entre a nossa alma e o corpo é \emph{pessoal}, como \emph{pessoal} é a união do \textsc{FiLHo} DE Dmgus com a natureza humana, e por isso, assim como todas as nossas operações se attribuem ao mesmo \emph{eu}, assim tambem 4 Pessoa do \textsc{FILHO} DE \textsc{DEus} se attribuem todas as operações, não só as da natureza divina, mas tambem as da natureza humana. --- Mas existe uma grande differenca; porque a união da nossa alma com o corpo é tambem \emph{natural}, pois em nós existe uma só \emph{natureza}, ao passo que a união do \textsc{FILHO} DE \textsc{Deus} com a natureza humana é só \emph{pessoal}, pois em N. S. Jesus \textsc{CHRISTO} as duas naturezas ficaram unidas, mas distinctas.}.

45. A união da \textbf{alma} com o corpo é substancial.

A união póde ser \emph{substancial} ou \emph{accidental}. --- Eº \emph{substancial}, quando varios elementos constituem, pela sua união, \emph{um só ser substancial}; p. e., \emph{um} só natureza, \emph{um} só supposto ou \pg{63}  pessoa. -- Eº \emph{accidental}, quando varios elementos constituem, pela sua união, \emph{um só ser accidental}; p. e. um \emph{aggregado} \footnote{A \emph{união}, geralmente, é a conjuneciüo de varios elementos para a formação de um ser composto. Se os elementos são \emph{redes}, a união é \emph{real}; se os elementos são \emph{ideaes} ou \emph{logicos}, a união é \emph{ideal} ou \emph{logica}.\par A união \emph{real} pode ser \emph{substancial} ou \emph{accidental}. --- E' \emph{substancial}, quando o composto, que resulta de varios elementos, é um todo \emph{substancial}; \emph{é accidental}, quando o composto é um todo \emph{accidental}. --- Desenvolvamos este ponto, que é da maxima importancia.\par São varios os modos da: união \emph{substancial}. --- Acontece, ás vezes, que dos dois elementos, que se unem, nenhum tenha subsistencia propria, isto é, nenhum exista por si e em si. Neste caso, esses dois elementos comecam a existir e existem somente no composto e pelo composto, e é o composto que propriamente subsiste em si. Assim a alma e \emph{o} corpo dos animaes irracionaes subsistem unicamente no composto. --- Pode tambem acontecer que um dos dois elementos tenha subsistencia propria, isto é, exista em si e por si, independentemente do outro elemento e do composto. Neste caso, o elemento subsistente communica a propria existencia, o proprio ser substancial ao outro elemento, e com este constitue o composto, que por isso subsiste pela subsistencia d' aquelle elemento mais nobre. --- Neste segundo caso, podem conceber-se duas hypotheses. Por quanto, o elemento subsistente, que communica ao outro elemento o proprio ser substancial, ou é \emph{incompleto} na sua especie, ou é \emph{completo}. Se o elemento \emph{subsistente} é \emph{incompleto} na sua \emph{especie}, refere-se ao outro elemento como ao complemento intrinseco da sua especie, e está para elle como a parte \emph{formal} está para a parte \emph{potencial}, como a \emph{forma substancial} esta para a materia \emph{prima}; e a união d' estes dois elementos é rigorosamente \emph{substancial}, porque d'ella deriva uma só natureza, um só supposto. Se, porém, o elemento subsistente é \emph{completo} na sua \emph{especie}, então esse ser subsistente, --- se a \emph{finito}, não pode attrahir a si um outro elemento e communicar-lhe a propria subsistencia, porque, nos entes finitos, a subsistencia é por tal modo proporcionada á essencia, que a esta se limita e nesta se exgotta; --- mas, se é \emph{infinito}, pode assumir e attrahir para a unidade do seu ser uma essencia finita, singular e completa, de modo que esta não tenha existencia e subsistencia propria, mas exista e subsista na existencia e subsistencia do Supposto infinito; e esta união é \emph{substancial}, porque, embora não importe unidade ou mistura de naturezas, importa todavia unidade de supposto, e um só supposto tem um só ser substancial, uma só subsistencia. Ensina-nos a Fé que esta união se realizou no ineffavel e dulcissimo Mysterio da Encarnação, no qual a Pessoa adoravel do Verbo divino assumiu a natureza humana, communicando-lhe o proprio ser substancial, a propria subsistencia divina, de modo que as duas naturezas --- a divina e a humana --- se encontraram unidas na unica Pessoa. do Unigenito Filho de Deus.\par A união \emph{accidental} é de duas especies. --- Podem varias substancias, completas e distinctas, unir-se por tal modo que cada uma conserve a pro- Dria subsistencia, e assim não ha entre ellas senão a unidade que se chama de \emph{aggregação}, --- unidade, que é apenas \emph{accidental}; tal é a união de\pgn{64}muitas pedras de um edificio. --- Pode dar-se tambem que a uma substancia singular e subsistente se unam diversas qualidades, que dão uma ulterior perfeição á propria substancia, já constituida no seu ser substancial; esta união tambem é \emph{accidental}; tal é o homem, que, permanecendo homem, começa a ser \emph{bom}, \emph{douto}, ecc. --- A segunda especie de união é menos imperfeita que a primeira.}.

Feita esta distineg\&o, prova-se a these.

\pg{64} A alma e o corpo constituem, pela sua união, \emph{uma} só pessoa e \emph{uma} só natureza. Ora a alma e o corpo não poderiam constituir, pela sua união, uma só natureza e uma só pessoa, se essa união não fosse \emph{substancial}; pois que, quando a união é \emph{accidental}, as naturezas e os suppostos são tantos, quantas são as substancias que se unem. Logo á união da alma com o corpo é \emph{substancial} \footnote{\emph{Substancia} significa não somente a \emph{natureza} ou \emph{essencia} de uma coisa, mias tambem o \emph{supposto} ou a \emph{pessoa}. Ora é nº essa duplice significação que a união da alma com o corpo se diz \emph{substancial}, emquanto d' ella resulta uma duplice unidade: de \emph{natureza} e de \emph{pessoa}. S. Thomaz diz: « Ex anima et corpore constituitur in unoquoque nostrum duplex unitas: \emph{naturae} et \emph{personae}. \emph{Naturae} quidem, secundum quod anima unitur corpori, \emph{formaliter} perficiens ipsum, ut ex duobus fiat una natura, sicut ex actu et potentia, vel materia et forma. Unitas vero \emph{personae} constituitur ex eis, in quantum est unus aliquis subsistens in carne et anima » (\emph{Sum}. \emph{theol}., 8 p., q. 2, a.1 ad 2). --- Note-se que esta proposição --- a \emph{naturesa humana é uma} só --- não differe da seguinte --- a \emph{essencia} ou \emph{substancia humana} é uma só. Por quanto, embora os conceitos de \emph{essencia}, de \emph{natureza} e de \emph{substancia} sejam formalmente diversos, comtudo no ente, que é \emph{substancia} (e não accidente), a \emph{essencia}, \& \emph{natureza} e a \emph{substancia} exprimem sempre uma e a mesma realidade; e por isso, havendo no homem uma só natureza, deve dizer-se que ha n'elle uma \emph{sb} substancia completa na propria especie.\par S. Thomaz trata d'esta questão da \emph{união da alma com o corpo} nas seguintes obras: --- \emph{Sum}. \emph{theol}., 1 p., q. 16; --- \emph{Contra gent}., lib. 2, ce. 56 e seg.; --- \emph{De spirit}. \emph{creaturis}, aa. 1e8; --- \emph{De anima}, aa. 1, 2, 9.\par a) Objectam: \emph{A} alma tem o seu ser \emph{completo}, independentemente do corpo. Ora a união de uma coisa com uma outra que tenha o seu ser \emph{completo é accidental}, e não \emph{substancial}. Logo a união do corpo com a alma não é \emph{substancial}, mas \emph{accidental}.\par Respondemos, com S. Thomaz (\emph{De anima}, a. 1 ad 1), que, embora a alma possua um ser \emph{completo}, comtudo não se segue que o corpo lhe se una por um modo \emph{accidental}, --- não só porque o ser que a alma communica ao corpo se torna o unico ser substancial do composto humano, e a união de duas coisas no mesmo ser substancial não pode deixar de ser \emph{substancial}, --- mas tambem porque se o ser da alma é \emph{completo} emquanto á \emph{subsistencia} (pois não precisa do corpo para subsistir), não é comtudo completo emquanto á \emph{especie}, e por isso, unindo-se o corpo á alma para a completar na especie, essa união é verdadeiramente \emph{substancial}.\par b) Continuam: Uma união \emph{substancial} da alma com o corpo é impossivel. Por quanto á alma, por ser \emph{simples} e \emph{indivisivel}, é heterogenea com\pgn{65}o corpo, que é \emph{extenso} e \emph{divisivel}. Ora o \emph{simples} não póde unir-se \emph{substancialmente} com o \emph{extenso}, nem o \emph{Indivisivel} com o \emph{divisivel}.\par Coneordamos em que o ente \emph{simples} e \emph{indivisível} é heterogeneo com o \emph{extenso} e o \emph{divisivel}; mas negamos que não possam estar \emph{substancialmente} unidos. Pelo contrario, o \emph{extenso} e o \emph{divisivel}, por isso mesmo que é \emph{extenso} e \emph{divisivel}, deve estar unido com o \emph{simples} e com o \emph{indivisivel}; visto que o \emph{extenso} e o \emph{divisivel} não póde adquirir nem conservar a unidade, se não por um principio que seja simples e indivisivel. Ora o nosso corpo é extenso e divisivel; para a sua formação concorrem \emph{quatorse} corpos \emph{simples}, innumeros principios immediatos e todos os elementos que selhe unem pelas funcções da nutrição. Mas, se o corpo, emquanto resulta de muitos elementos, é multiplo; todavia, emquanto é vivo, é substancialmente \emph{uno}, pois em todo o nosso organismo existe uma perfeita harmonia, e todas as nossas operações concorrem para um unico fim, que é o bem e a conservação do individuo. Esta \emph{unidade substancial} deve derivar de um principio \emph{simples} e \emph{indivisível}; pois a unidade não póde derivar de um principio \emph{composto} e \emph{divisivel}, aliás o effeito seria superior á causa. Este principio \emph{simples} e \emph{indivisivel} deve ser \emph{intrinseco} ao proprio corpo; pois um principio \emph{extrinseco} não póde communicar a unidade \emph{substancial} a seres varios e diversos. Este principio \emph{simples}, \emph{indivisivel} e \emph{intrinseco} ao nosso corpo é a alma. Logo o corpo não só não exclue, mas exige a sua união com a alma, como factor indispensavel da sua \emph{unidade} substancial. --- Tambem, para a nossa alma, é mais natural a sua união com o corpo, do que com uma outra alma; pois é mais natural a união de um acto com uma potencia do que com outro acto, pois dois actos formam um \emph{aggregado}, ao passo que o acto e a potencia formam um todo completo.\par 0) Insistem: O espirito e a materia sao substancias contrarias. Ora substancias contrarias excluem-se e não se podem unir \emph{substancialmente}.\par Resposta. Repugna que duas substancias contrarias se possum attribuir a0 mesmo sujeito, ao mesmo tempo e sob o mesmo aspecto, porque neste caso excluem-se mutuamente; mas não repugna que duas coisas, contrarias entre si, possam unir-se para constituir uma terceira, que não seja nem uma nem outra, mas um composto de ambas. Ora a alma e o corpo, por serem substancias contrarias entre si, não podem attribuir-se ao mesmo sujeito, ao mesmo tempo e sob o mesmo aspecto, porque é absurdo que um ente, que é espirito, seja tambem corpo; mas não ha nenhuma contradicção em que á alma e o corpo se unam para formar uma \emph{terceira substancia} --- o homem, que não é nem corpo só, nem alma só, mas um composto de alma e de corpo.\par á) Replicam: Entre duas coisas, que se unem \emph{substancialmente}, deve existir uma relação. Ora entre a alma intellectual e o corpo não existe nenhuma relação. Logo a alma não póde unir-se \emph{substancialmente} ao corpo.\par Resposta. Concedemos que entre duas coisas, que se unem \emph{substancialmente}, deve existir uma certa relação; mas não é necessario que esta relação seja de \emph{natureza}, emquanto as duas substancias devam possuir á mesma natureza, mas basta que seja uma relação de \emph{proporção}, emquanto uma se refira a outra, como o \emph{acto á potencia}. Ora, se entre a alma intellectual\par 5 Vor. II.\pgn{66}e o corpo não existe uma relação de \emph{natureza}, existe uma relação de \emph{proporção}; pois, como provaremos, a alma, pelo facto de vivificar o corpo e de o constituir numa certa e determinada especie, é \emph{acto} do mesmo corpo.}. :

\pg{65} 46. A \textbf{alma} humana \textbf{esta} unida ao corpo, como a \textbf{forma substancial} á \textbf{materia}.

a) A alma e o corpo, pela sua união, não poderiam constituir, como constituem, uma só \emph{pessoa} e uma só \emph{natureza}, se \pg{66}  á alma não se unisse ao corpo, como a forma \emph{substancial} á materia. Na verdade, para que dois elementos formem um só supposto e uma só natureza, é necessario que um .d'elles seja activo, proprio e determinante, e o outro seja passivo, commum e determinavel. Ora o elemento activo, proprio e determinante, diz-se \emph{forma substancial}; \emph{e} o elemento passivo, commum e determinavel, diz-se \emph{materia}. Logo a alma está unida ao corpo, como a forma \emph{substancial á materia} \footnote{Para a exacta comprehensão d'esta these, deve-se ter presente a theoria ácerca da \emph{composição essencial dos corpos} (Oosmologia, cap. III). Todos os corpos são essencialmente compostos de dois principios: da \emph{materia prima} e da \emph{forma substancial}. \emph{A materia prima} é uma substancia incompleta e determinavel, não tem forma nem especie, mas pode receber toda e qualquer forma e pode ser determinada a toda e qualquer especie de substancias corporeas. A \emph{forma substancial} é uma substancia incompleta e determinante, que, unindo-se com a materia communicadlhe o ser \emph{primeiro}, que é o ser \emph{substancial}, e assim a determina e colloca nº uma certa especie, de modo que seja tal ente, e não outro. --- E, como \emph{ente} e \emph{uno} são, na realidade, uma e a mesma coisa, segue-se que a forma, a qual communica á materia a \emph{entidade}, communica-Ihe tambem a \emph{unidade}.\par O argumento, tirado da unidade da nossa \emph{pessoa} e da nossa natureza, é convincente. Na verdade, se a alma e o corpo fossem duas realidades completas e dotadas do seu acto \emph{especifico}, constituiriam dois suppostos e duas naturezas distinctas, e não participariam do mesmo ser substancial. Logo, para que a alma e o corpo constituam uma só pessoa e uma só natureza, e participem do mesmo ser substancial, é necessario que uma d'essas duas substancias communique o ser e a outra o receba, que uma seja o principio de especificação e a outra seja o sujeito determinado a uma certa especie de seres finitos. A substancia, que communica o ser e a especie, é a \emph{alma}; e a substancia, que recebe o ser e a especie, é o \emph{corpo}; pois é pela alma que o nosso corpo é humano, vegeta e sente. A alma, pois, é forma \emph{substancial} do corpo.\par Esta conclusão é certissima. Por que modo poderiam unir-se a alma e o corpo? --- Por simples \emph{mistura} ? Não; porque, nas misturas, as particulas de dois corpos confundem-se, mas os dois corpos conservam a propria natureza e não constituem uma substancia nova; ora a alma humana não tem partes, que se possam confundir com as partes do corpo, e a alma e o corpo constituem, pela sua união, uma nova substancia completa. --- Por \emph{combinação chimica}? Tambem não; porque a combinação chimica só é possivel entre dois corpos, que por mutua acção se alteram e se modificam essancial e substancialmente, de modo que resulta uma nova forma, um novo ser, que não existia antes da conbinação; mas á alma humana, sendo simples e espiritual, é absolutamente incorruptivel, não pode estar sujeita a transformações, não pode perder o proprio ser nem confundir-se com outros seres, quer materiaes quer espirituaes. --- Por \emph{contacto}? Tambem não; porque, embora não repugne que uma substancia espiritual possa\pgn{67}\textbf{anover} o corpo, todavia tal união não seria \emph{substancial}, como effectivamente é, mas seria \emph{accidental}, como é a união de duas substancias que apenas se tocam. Logo --- ou deve negarse que a união da alma com o corpo é \emph{substancial}, e, nº esse caso, deve tambem megar-se que a alma e o corpo constituem, na sug união, um só individuo humano, --- ou deve concluir-se que a alma se une ao corpo, como a \emph{forma substancial á materia prima}.}.

\pg{67} ---b) Forma \emph{substancial} é a realidade que communica a um \textbf{ente} o ser \emph{substancial}, de modo que o determina a uma certa especie, e com elle constitue uma só substancia completa. Ora “a alma \emph{Humana} communiea ao corpo o \emph{ser humano}, que é \emph{substancial}, de modo que o determina á especie dos corpos \emph{humanos}, e com elle constitue uma só substancia completa, --- o \emph{homem}. Logo a alma é forma \emph{substancial} do corpo, isto é, está unida ao corpo como a forma \emph{substancial} á materia \footnote{S. Thomaz expõe este mesmo argumento. Para que, diz o Angelico Doutor, uma coisa seja \emph{forma substancial} de outra, são necessarias e sufficientes duas condições. A primeira condição é que a tal coisa communique, pela sua união, o ser \emph{substancial}. a outra, de que se diz forma, e assim a determine e colloque nº uma certa especie de substancias corporeas. A segunda condição, consequencia da primeira, ó que a mesma coisa constitua com a outra, de que é forma, um só ser substancial, que é o ser da propria substancia composta. Ora estas duas condições verificam-se no respeito da alma ao corpo. Na verdade, a alma constitue o corpo e o colloca na especie dos entes \emph{vivos}, e por isso communica-lhe o \emph{ser substancial}; porque, nos entes vivos, a vida não é uma qualidade accidental, mas é a propria substancia, identifica-se com o proprio ser substancial: \emph{vivere viventibus est esse}. Além d'isso, communicando-se ao corpo, a alma eleva-o para si, para o proprio ser substancial, que ella recebeu de Deus por creação e possue independentemente da sua união com o corpo; de modo que o corpo participa do proprio ser substancial da alma, e assim um e o mesmo é o ser da alma e do corpo, isto ó, do composto humano. Logo a alma é \emph{forma substancial} do corpo. (\emph{Contra Gentes}, l. Il, c. 68). --- Em poucas palavras: A alma communica ao corpo a vida. Ora a vida é o proprio ser substancial do vivente. Logo a alma communica ao corpo o ser substancial. Mas o que communica o ser substancial é e diz-se \emph{forma substancial}. Logo a alma é forma substancial do corpo (Cf. S. Thom., \emph{De anima}, a. 1).\par \emph{Um e o mesmo é o ser da alma e do corpo}. Dilucidemos esta proposição, que tem um immenso alcance. A alma humana é, ao mesmo tempo, forma \emph{substancial} e \emph{subsistente}, --- Emquanto forma \emph{substancial}, a alma une-se pela sua essencia ao corpo e, n'essa união, communica-lhe o ser substancial e 2 especie. « Totum corpus, diz S. Thomaz, et omnes ejus partes habent \emph{esse substantiale} et \emph{spectficum} per animam » (\emph{De spirit}. \emph{creat}., a. 4. --- Emquanto - \emph{subsistente}, a alma tem um ser proprio, independentemente do corpo, e este - Ser é porella communicado ao corpo, e assim o corpo, que não tem ser proprio, é elevado á participação do ser da alma, de modo que um e o mesmo é o ser da alma e do corpo. « Anima, continua o mesmo S. Doutor, habet \textbf{esse} subsistens in quantum esse suum non dependet a corpore, utpote supra - materiam corporalem elevatum; et tamen ad hujus esse communionem\pgn{68}recipit corpus, ut sic sit \emph{unum esse} animae et corporis, quod est esse hominis « (\emph{De spirit}. \emph{creat}., a. 2, ad 8). --- 'Unindose, porém, ao corpo emquanto Eo \emph{essencia}, a alma não lhe se une emquanto á \emph{totalidade} da \emph{sua energia} (como o demonstram as operações intellectuaes) Por isso o ser da alma não está comprehendido ou immerso inteiramente na materia, ma é apenas \emph{attingido} por ella, e permanece subsistente, e na propria alma reside como no unico sujeito capaz de o receber em toda a sua plenitude. Diz S. Thomaz: « Anima humana propter suam nobilitatem supergreditur facultatem materiae corporalis, e\& non potest totaliter includi ab ea. Unde remanet aliqua actio, in qua materia corporalis non communicat ». (\emph{De anima}, YII, 1. 7).\par Objectam: E' impossivel que um e o mesmo seja o ser da alma e do corpo; por quanto o ser do corpo é composto e corruptivel, ao passo que o ser da alma é simples e incorruptivel.\par Resposta. A difficuldade confunde o ser da \emph{essencia} ou natureza com. o ser da \emph{existencia}; porque o \emph{ser} pode significar tanto a \emph{essencia} ou \emph{natureza}, como a \emph{existencia}. Se por \emph{ser} se entende a \emph{essencia} ou \emph{natureza}, então é evidente que não é um e o mesmo o ser da alma e o ser do corpo, porque á .essencia da alma é simples e incorruptivel, ao passo que a essencia do corpo é composta e corruptivel; e por isso temos necessariamente duas essencias, porque temos duas substancias. --- Ma se por \emph{ser} se entende a \emph{existencia}, como no. caso presente, então um eo mesmo é o ser da alma e do corpo; porque o corpo não tem existencia propria, mas recebe-a da alma, de modo que a alma e o corpo subsistem pela unica existencia da alma. --- Esta unica \emph{existencia}, sendo a ultima actualidade da alma humana e, por meio d'esta, do corpo, é um \emph{acto} e um acto \emph{subsistente}. Por ser \emph{acto}, a \emph{existencia} ou o ser da alma humana, e por isso do corpo, é \emph{sz»iples}, e não se corrompe \emph{por si}, como se corrompe o ente composto, --- e por ser acto \emph{subsistente}, o ser da nossa alma, e por isso do corpo, não se corrompe pela corrupção do sujeito em que é recebido, como se corrompe a alma dos brutos. Portanto o \emph{ser} do homem é sempre simples e incorruptivel, quer se considere na alma como no seu principio e sujeito proporcionado, quer se considere no corpo, no qual é recebido como em sujeito participante e incapaz de receber toda a plenitude da perfeição, de que o mesmo ser é dotado. E quando se corrompe a natureza do homem, não se corrompe o seu ser ou a sua \emph{existencia}, porque o ser é proprio da alma, e, como a alma é espiritual, incorruptivel e immortal, o ser continua a permanecer nº ella, por toda a eternidade (Cf. \emph{Sum}. \emph{theol}., Y p., q. 76, a. 1 ad 5).}.

c) A alma ou é forma \emph{accidental}? do nosso corpo, ou é forma \emph{substancial}, ou é uma substancia que lhe \emph{assiste}. Não póde dar-se outra hypothese. Ora a alma não é forma \emph{accidental} do nosso corpo; pois a forma \emph{accidental} communiea uma perfeição \emph{accidental}, e a nossa alma, que é \emph{substancia}, embora \emph{incompleta}, communica ao corpo o ser \emph{substancial}. Nem é uma substancia que \emph{assiste} ao corpo; porque, nesse caso, a união seria apenas \emph{accidental}, pois a alma e o corpo não constituiriam uma só substancia \emph{completa}, mas um \emph{aggregado} de substancias, e a união que existe no \emph{aggregado É accidental}. Logo \pg{68} \pg{69}  é forma \emph{substancial} do corpo, porque é \emph{substancia}, \emph{e} commu- \textasciicircum{}iea ao corpo o ser \emph{substancial}. Logo a alma está unida ao --- corpo, como a forma \emph{substancial} á materia \footnote{Dizem: A alma humana é substancia espiritual e subsistente. Ora uma substancia espiritual e subsistente nao pode ser acto ou forma do corpo. Logo a alma não pode ser forma substancial do corpo.\par Já demos os principios, que resolvem esta difficuldade. A alma humana !é substancia espiritual e subsistente, mas é incompleta na sua especie, E porque sem o concurso do corpo não pode absolutamente exercer as operações da vida sensitiva e vegetativa, nem, no estado presente, as da vida intellectual. Incompleta na sua especie, a alma tem uma intrinseca e essencial relação ao corpo, e só no corpo e com o corpo ella se torna completa na sua especie, porque só com o corpo e no corpo pode exercer \emph{todas} as operações proprias do homem. --- Ha, na verdade, uma triplice especie de formas. Algumas são subsistentes e completas na sua especie, e dizem-se formas \emph{separadas}, non podem ser actos da materia: como são os Anjos. Outras formas não subsistem, não existem em si, não têm ser proprio, mas têm o ser no composto, e n'este existem e subsistem, de modo que não podem existir fóra do composto, e portanto dependem da materia no proprio ser, e dizem-se formas \emph{materiaes}, ou formas \emph{informantes}, mas \emph{não subsistentes} ; taes são as almas dos brutos. Outras, intermedias entre as formas \emph{separadas} e as materiaes, têm o seu ser independente da materia e por isso, separadas da materia, perseveram na existencia; mas, como o seu ser não é completo na especie, são destinadas a unir-se substancialmente com a materia, e com ella e nella se completam na especie, e estas dizem-se \emph{informantes} e \emph{subsistentes}: tal é a alma humana. Por isso a nossa alma, unindo-se ao corpo como forma \emph{substancial}, e constituindo com elle uma especie completa, aperfeicoa-se no seu \emph{ser} (que se torna \emph{completo} na \emph{especie}) e tambem nas suas \emph{operações}, que exigem o concurso do corpo. (Cf. S. 'Thom., \emph{De spzrit}. \emph{creat}., a. 2).\par b) Continuam: A forma, que dà o ser á materia, é \emph{material}. Ora a alma humana não é \emph{material}, mas é \emph{espiritual}. Logo a alma humana não dá o ser á materia, e por isso não é forma \emph{substancial} do corpo.\par A solução está contida na doutrina exposta. A forma, que dá o ser á materia, é \emph{material}, isto é, é dependente da materia na existencia, quando 0 seu ser é totalmente comprehendido pela materia e immerso na materia; mas não se pode dizer nem é \emph{material}, quando o seu ser é apenas \emph{atzingido} pela materia, de modo que excede sempre a capacidade da materia. Ora o ser da alma, embora seja communicado á materia, não é inteiramente comprehendido por ella, mas excede immensamente a sua capacidade (como o mosíram as operações intellectuaes); e por isso a alma não é e não se pode dizer \emph{material}, mas simplesmente \emph{espiritual} (Cf. \emph{Contra Gentes}, 1. 2, c. 58).\par 0) Accrescentam : A intelligencia humana não pode ser acto ou forma do corpo, porque é espiritual. Ora a alma não é menos espiritual do que uma sua faculdade. Logo a alma não pode ser acto ou forma do corpo.\par Resposta. A. intelligenecia não pode ser acto ou forma de um orgão corporeo, não só porque é principio proximo de uma operação immaterial, da qual não pode ser comprineipio um orgão corporeo, mas tambem porque é forma \emph{accidental}. Na verdade, á forma \emph{accidental} recebe o ser, e recebe-o do\pgn{70}sujeito, de que é acto e que é \emph{em acto}, \emph{e} no seu ser depende do mesmo sujeito; e por isso se a nossa intelligencia fosse acto ou forma de um orgão corporeo, dependeria, no seu ser, de um sujeito \emph{material}, e seria forma \emph{material}, e não \emph{espiritual}. Pelo contrario, a forma \emph{substancial} e \emph{subsistente}, qual é a nossa alma, não recebe o ser do sujeito, a que está unida, mas é ella que communica o proprio ser ao sujeito, que por si é só \emph{em potencia}; e por isso não é a forma que depende, no seu ser, do sujeito, mas antes é o sujeito que depende da forma, e esta, na união, não se torna \emph{material}, mas permanece sempre subsistente e \emph{espiritual}. Logo, ainda que a espiritualidade da alma não seja inferior á espiritualidade da intelligencia, não admira que a alma possa ser forma do corpo, e não o possa ser a intelligencia. --- D'onde se vê que a nossa alma é forma do corpo quanto á \emph{essencia}, e não quanto á faculdade \emph{intellectual}. (Cf. \emph{Sum Theol}., I p. q. 76, a 6).\par Pergunta-se: A alma é forma do corpo emquanto é \emph{intellectual}?\par 1 A pergunta póde tomar-se em dois differentes sentidos; por quanto a palavra « \emph{intellectual} » póde referir-se --- ou á propria \emph{essencia} da alma, ---- ou ao mais elevado grau da sua perfeição. --- Se a palavra « \emph{intellectual} n se refere á propria essencia da alma, então deve dizer-se que a alma se une ao corpo, emquanto é \emph{intellectual}; pois é essencialmente \emph{intellectual}. --- Se á palavra « \emph{intellectual} » se refere ao mais elevado grau da sua perfeição, respondemos que a alma se une ao corpo, não emquanto é \emph{intellectual}, mas emquanto é \emph{virtualmente} vegetativa e sensitiva. --- Repetimos: A alma une-se ao corpo, quanto á sua \emph{essencia}, e não quanto ao grau de vida \emph{intellectual}, que a materia é incapaz de receber.\par \emph{LE}\par *\par Os philosophos, tanto antigos como modernos, apresentaram muitas e, diversas soluções do problema ácerca da união da alma com o corpo. Limitamo-nos a resumir e criticar as mais celebres, que são as de \emph{Platão}, \emph{Mallebranche}, \emph{Leibnitz}, \emph{Leclerc}, \emph{Locke} e \emph{Rosmini}.\par A) \emph{Platão} ensinou que a alma esta unida ao corpo, como o motor ao movel, o musico á lyra. D'hai concluiu que só a \emph{alma} constitue o homem.\par Esta opinião é falsa pelos seguintes fundamentos: \emph{P}\par a) Destroe a unidade \emph{substancial} do homem. Na verde, a alma e o corpo formam, na sua união, \emph{uma só substancia completa}. Ora ninguem dirá que o motore o movel, o musico e a lyra, o artifice e o instrumento, formam, na sua união, uma só \emph{substancia completa}.\par 5) O nosso corpo desums da alma a especie de corpo \emph{umano}. Ora não poderia desumir da alma a especie, se a alma se lhe unisse como o motor ao movel; pois e elemento, de que um ente desume a sua especie, deve constituir a essencia do proprio ente, mas o motor \emph{está fóra} do movel.\par B) \emph{Mallebranche}, seguindo as doutrinas de \emph{Descartes}, suppõe a alma e o corpo destituidos de actividade, por fórma que nem a alma opéra no corpo, nem este na alma. Se, recebidas algumas impressões no corpo, se suscitam na alma conhecimentos e appetites correspondentes; e se, posto o acto da vontade, o corpo se determina aos relativos movimentos, isto deve attribuir-se\pgn{71}á acção immediata de Deus, que, segundo este philosopho, é a unica causa efficiente, a que devem referir-se todas as operações do homem. Nesta correspondencia de movimentos e operações consiste a união da alma com o corpo. --- E o systema das causas \emph{occasionaes}.\par O systema das \emph{causas occasionaes} deve rejeitar-se pelas seguintes razões:\par \emph{a}) \emph{Baséa-se num falso fundamento}. Com effeito, \emph{Mallebranche} parte do principio de que as creaturas são destituidas de toda a actividade propria e intrinseca. Ora este principio é falso; porque, como demonstrámos, não só não repugna que as creaturas possuam uma actividade propria e intrinseca, mas effectivamente a possuem.\par \emph{b}) \emph{Segue un falso methodo}. Na verdade, \emph{Mallebranche}, para explicar as relações entre a alma e o corpo, recorre immediatamente á Causa primeira. Ora este methodo não é digno de um philosopho, que, presupposta a existencia de Deus, como Causa \emph{primeira} e \emph{universal} de todos os phenomenos da natureza, deve procurar as causas \emph{segundas} e \emph{particulares}.\par \emph{c}) \emph{Destroe a unidade substancial do homem}. Por quanto, se a alma e o corpo, na sua união, constituissem uma só substancia, pelo facto de que Deus, na occasio de um movimento do corpo, produz na alma uma affeição correspondente, e vice-versa; tambem formariam uma só substancia duas machinas collooadas uma perto de outra, e movidas pelo artifice de modo que os moyimentos de uma correspondessem aos movimentos da outra.\par \emph{d}) \emph{Leva a perniciosas consequencias}. Por quanto, se Deus fosse a \emph{unica} causa dos phenomenos que se realizam em nós, a Deus deveriamos attribuir \emph{todas} as nossas cogitações e operações, não só as bóas e honestas, mas tambem as criminosas e más. Ora isto repugna á infinita perfeição de Deus, e destroe todo o fundamento da moralidade. :\par €) \emph{Leibnits} diz que a alma e o corpo são duas substancias, duas monadas, das quaes uma não opéra na outra, porque a acção \emph{transeunte} repugna, mas que se desenvolvem numa serie de operações, das quaes a seguinte tem a sua razão sufficiente na anterior, sem nenhuma dependencia do corpo para com a alma, nem d'esta para com aquelle. Se as operações da alma correspondem ás do corpo, e vice-versa, isto succede porque, entre infinitos espiritos e infinitos corpos, Deus escolheu e uniu aquelle espirito e aquelle corpo, cujas operações estivessem em perfeita harmonia. Nesta harmonia consiste a união da alma com o corpo. --- Eo systema da \emph{harmonia preestabelecida}.\par Este systema não póde acceitar-se pelas seguintes razões:\par \emph{a}) \emph{Baséa-se num falso fundamento}. Por quanto, \emph{Leibnits} nega a acção da alma no corpo e d'este naquella, porque diz que a acção \emph{transeunte} repugna. Já demonstrámos que a acção \emph{transeunte} não sô não repugna, mas effectivamente se dá nas creaturas.\par \emph{b}) \emph{Destroe a unidade substancial do homem}. Na verdade, se a alma e o corpo constituissem uma só substancia, porque Deus move o corpo em harmonia com os movimentos da alma, então dois relogios formariam uma só substancia pelo facto de serem movidos pelo mesmo artifice para marcar 2 mesma hora, o que é absurdo.\par \emph{c}) \emph{Tira a liberdade ao homem}, Com effeito, este systema sustenta --- que\pgn{72}as operações do corpo, na sua producção, obedecem a leis mechanicas, e a seguinte tem a sua razão sufficiente na anterior, --- que as operações da alma devem corresponder \emph{necessariamente} ás operações do corpo, e que á seguinte operação é \emph{determinada} pela anterior; e assim Deus seria a causa de todas as nossas obras, boas e más; o que é absurdo e impio.\par \emph{d}) \emph{Leva ao idealismo}. Por quanto, se a percepção, que temos dos corpos, não fosse determinada pela acção dos proprios corpos nas faculdades perceptivas, não haveria razão sufficiente para affirmar que os corpos existem.\par D) \emph{J}. \emph{Leclerc} disse que a alma está unida ao corpo por meio de uma terceira natureza, que participa da alma e do corpo, e produz toda a harmonia que existe entre os pensamentos da alma e as operações do corpo; pois que essa natureza, embora não possua a consciencia dos seus actos, todavia refere á alma o que se passa no corpo, e executa no corpo as operações que a alma prescreve. --- E' o systema do \emph{mediador plastico}.\par Este systema é falso pelas seguintes razões:\par a) Destroe a unidade \emph{substancial} entre a alma e o corpo; pois deve admittir que estas duas substancias, pelo facto de operarem uma independentemente da outra, constituem duas \emph{naturezas} e dois \emph{suppostos}, mas nunca uma \emph{substancia completa}, como é o homem.\par \emph{b}) Se entre a alma e o corpo existisse uma terceira substancia, esta deveria possuir, ao mesmo tempo e sob o mesmo aspecto, todos os attributos da alma e todos os attributos do corpo, e por isso deveria ser simples e extensa, corporea e espiritual, racional e irracional, ete. Ora repugna que uma e a mesma entidade seja dotada de propriedades contrarias.\par E) \emph{Locke}, seguido por \emph{Newton}, \emph{Clarke}, \emph{Macke}, \emph{Storchenauw}, opina que a alma e o corpo são duas substancias \emph{completas}, dotadas de operações proprias e independentes. A alma immediatamente influe e opera no corpo, emquanto o determina ao movimento; e, vice-versa, o corpo influe e opéra na alma, emquanto os movimentos do nosso corpo, determinados pela acção dos corpos circumstantes e transmittidos ao cerebro, determinam a alma á percepção da existencia e das propriedades dos corpos. Neste mutuo influxo da alma no corpo, e d'este naquella, consiste a união das duas substancias. --- E' o systema do \emph{influxo physico}.\par O systema do \emph{influxo physico} é falso pelas seguintes razões:\par a) \emph{Toma o effeito pela causa}. \emph{Locke} diz que a alma e o corpo se unem, emquanto exercem entre si um mutuo influxo. Ora isto é falso; pois esse mutuo influxo não é \emph{causa} da união, mas é \emph{e/feito} da propria união.\par \emph{b}) \emph{Destroe} a \emph{unido substancial da alma com o corpo}. Na verdade, a alma e o corpo sendo, na opinião de \emph{Locke}, duas substancias \emph{completas}, só poderão formar um \emph{aggregado} de substancias; e por isso a sua união será apenas \emph{accidental}, como o é toda a união que se funda nas operações.\par c) \emph{Favorece o materialismo}. Com effeito, o corpo é extenso e material. Ora um agente material só póde operar num ser material mas nunca num ser immaterial ou espiritual. Logo, se se admittir que o corpo opéra na alma, deverá tambem admittirse que a alma é material.\par F) \emph{Rosinini} propos a seguinte solução: A substancia da alma é constituida\pgn{73} por uma \emph{sensação}, que é anterior á consciencia e que se chama \emph{sensação fundamental}. A alma humana é \emph{sensitiva} e \emph{intellectual}, Emquanto \emph{sensitiva}, \emph{a} alma une-se com o corpo por meio da \emph{sensação} fundamental; pois \emph{sentir} o corpo é o mesmo que unir-se-]he e animal-o. Emquanto \emph{intellectual}, a alma une-se com o corpo por meio da percepção immanente da sensação fundamental. Por isso o nexo, que liga num individuo a alma racional e o corpo humano, consiste na primitiva, natural e continua \emph{percepção} da sensação fundamental e animal. --- D'onde se vé que \emph{Rosmini} colloca a união da alma com o corpo em duas percepções, --- uma \emph{sensitiva}, pela qual a alma \emph{sente} o corpo, --- outra \emph{intellectual}, pela qual a alma percebe a \emph{sensação f'undamental} (Cf. Rosmini, \emph{Psicologia}; --- Pestalozza, \emph{La mente di} A. \emph{Rosmini}).\par O systema de \emph{Rosmini} ácerca da união da alma com o corpo é falso --- a) nos seus \emph{fundamentos}, --- b) na sua \emph{essencia}, --- c) nas suas \emph{consequencias}.\par a) E' falso nos seus fundamentos. Por quanto:\par \emph{a}) \emph{Rosmini} colloca a essencia da alma numa \emph{sensação fundamental}, originaria e estavel, que é o unico principio ou sujeito de todas as outras sensações e operações humanas. Ora é falso que a essencia da alma consiste numa sensação. A alma humana é uma \emph{substancia}, subsiste em si e não precisa de sujeito a que adhira; ao passo que a sensação é \emph{accidente}, porque não subsiste em si, mas precisa de um sujeito a que adhira.\par \emph{b}) \emph{Rosmini} affirma que é só a alma que sente. Esta affirmacao é inexacta; porque as qualidades \emph{sensiveis} do corpo, que são o objecto \emph{proprio} da sensação, não podem impressionar uma substancia espiritual, como é á nossa alma, nem podem ser por ella percebidas.\par b) E' falso na sua \emph{essencia}.\par \emph{a}) \emph{Rosmini} faz consistir a união da alma com o corpo numa percepção ou sensação fundamental, que a alma tem do corpo. Ora isto é falso; porque, então, a união da alma com o corpo seria apenas \emph{accidental} e \emph{intencional}. --- Seria \emph{accidental}, pois toda a percepção ou sensação não é substancia mas accidente, e uma união que se basêa no \emph{accidente} não póde deixar de ser \emph{accidental}. --- Seria \emph{intencional}; visto que a união, que existe entre uma faculdade perceptiva e o objecto percebido, não é \emph{real}, pois o objecto não se une immediatamente na sua entidade á faculdade cognitiva, mas é apenas \emph{intencional}, emquanto o objecto se une á faculdade por meio de uma imagem ou similhança que o representa.\par \emph{b}) \emph{Rosmini} admitte que o corpo opéra na alma; pois diz que a alma não pôde \emph{sentir} o corpo, sem que seja determinada por elle ao acto da sensação (\emph{Psicol}, vol. I, lib. 9, c. v, a. 8). --- Esta esserção é falsa, porque, como demonstrámos contra \emph{o} systema do \emph{influxo physico}, o corpo nao póde operar numa substancia espiritual, como é a nossa alma.\par c) \emph{Rosmini} por um lado affirma que o corpo só opéra na alma em virtude da força, que a alma lhe communica no momento em que o sente, e por outro lado diz que a alma não póde sentir o corpo e por isso unir-se-]he, sem que o corpo a determine ao acto da sensação (\emph{Psicol}, vol. l, lib. 2, c. IX, a. 4). --- Estas affirmações contêm um \emph{circulo} vicioso; pois o corpo receberia da alma a força de operar só no momento em.que é \emph{sentido} por\par n\pgn{74}ella, e todavia deveria possuir tal força antes d'esse momento para poder determinar a alma ao acto da sensação.\par 0) E' falso nas suas \emph{consequencias}.\par a) \emph{Rosmini} admitte que, assim como a alma é forma substancial do corpo, assim tambem o \emph{corpo é forma substancial} da \emph{alma}; visto que a alma e o corpo actuam mutuamente uma noutro. (\emph{Psicol}, vol. I, lib. 2, c. IX, a. 2). --- Esta opinião é falsa. \emph{Forma substancial} é a realidade que communica á materia a especie e as propriedades. Ora não é o corpo que communica á alma a especie humana e a vida, aliás todos os corpos seriam humanos e seriam vivos e principios de vida; mas é a alma que constitue o corpo numa certa e determinada especie, e que lhe communica a vida.\par \emph{b}) \emph{Rosmini}, admittindo que a alma se une ao corpo \emph{por meio da sensação}, deve admittir que a alma é forma \emph{mediata} do corpo. Ora isto é falso. A alma une-se ao corpo por si mesma, communicando-lhe a propria entidade, e por isso é forma \emph{immediata} do corpo.\par \emph{c}) \emph{Rosmini}, partindo do principio de que não é a \emph{essencia} da alma que \emph{determina} o corpo, mas sim um producto d'ella, isto é, a \emph{percepção}, conclue que a alma não é a forma \emph{substancial} do corpo, mas é a causa \emph{efficiente} d" essa, forma (\emph{Psicol}. , vol. IL, lib. 1) --- Esta conclusão é tão falsa, como é falso o principio de que deriya. A percepção, quer seja sensitiva quer intellectual, nao póde ser forma \emph{substancial}, porque é apenas um \emph{accidente}.}.

47. \textbf{Entre} a alma e o corpo dá-se um mutuo influxo. --- Da \pg{70}  união da alma com o corpo resultam no homem muitas e diversas faculdades, --- as vegetativas, as sensitivas e as intellectuaes. Ora todas estas faculdades derivam da alma, como \pg{71}  \textgreater{}. * * * L , * do seu \emph{principio} commum; pois a alma é forma \emph{substancial} do corpo, e todas as forças, de que um ser é dotado, derivam da sua forma \emph{substancial}. Se a alma é o principio commum de \pg{72}  todas as faculdades, é claro que a actividade de uma faculdade deve terminar sempre na essencia da alma e d'esta refluir para as outras faculdades; e assim as faculdades superiores, que \pg{73}  são as \emph{intellectuaes}, exercerão um certo influxo nas faculdades 4 * J . = * pt inferiores, que são as \emph{vegetativas} e as \emph{sensitivas}, e estas naquellas. Ora as faculdades superiores residem, como em seu \pg{74}  \emph{sujeito}, na alma; e as faculdades inferiores residem no corpo informado pela alma. Logo entre a alma e o corpo deve existir e existe um. mutuo influxo \footnote{A these apoia.se na experiencia. Entre as faculdades organicas e as intellectuaes existe um mutuo influxo, uma dependencia mutua. As faculdades organicas, como vimos, influem nas intellectuaes; porquanto as faculdades sensitivas apresentam á intelligencia os materiaes do conhecimento, e por isso, se não estiverem no seu estado normal, a inteligencia não poderá perceber o objecto, como é em si mesmo, e o conhecimento será deficiente ou falso. --- As faculdades intellectuaes exercem uma reacção sobre as organicas; assim o trabalho excessivo do espirito impede o regular exercicio das faculdades inferiores e chega a produzir fortes commoções nervosas e até a alienação mental (Cf. S. Th. \emph{De verit}., q. 26, \emph{de pass}. \emph{animae}, a. 10).\par Ora estes phenomenos só podem explicar-se no systema, que admitte uma união \emph{substancial} entre a alma e o corpo. Porquanto, \& alma e o corpo, se fossem duas substancias \emph{completas} e não constituissem uma só natureza, não poderiam de modo algum exercer entre si o reciproco influxo que exercem; pois os movimentos do corpo e os pensamentos da alma, por serem operações \emph{immanentes}, levem ter o seu termo na propria substancia em que têm começo, e não numa substancia distincta.\par Note-se a differenca entre a nossa doutrina ácerca do \emph{mutuo influxo} entre a alma e o corpo e a doutrina dos defensores do \emph{influxo physico}. --- 1) Nós.admittimos o \emph{influxo plujsico}, mas como um \emph{effeito} e uma \emph{consequencia} da união da alma com o corpo; ao passo que \emph{Locke} faz consistir naquelle influxo a união das duas substancias. --- 2) Partindo do principio de que a alma é forma \emph{substancial} do corpo, sustentamos que o corpo opéra na alma, não emquanto é corpo, mas emquanto as faculdades organicas, que residem\pgn{75}no \emph{corpo}, como no seu \emph{sujeito}, terminando na essencia da alma, têm um reflexo nas faculdades \emph{intellectuaes}, que residem na propria \emph{alma}, como em seu \emph{sujeito}; ao passo que \emph{Locke}, por sustentar que o corpo não constitue com a alma uma substancia completa, conclue que o corpo, \emph{emquanto corpo}, opéra na alma, o que é impossivel, pois que a materia não póde operar no espirito.}.

\pg{75} 48. \textbf{Consequencias da} união \textbf{da alma} com o corpo. --- As \textbf{principaes} consequencias da união da alma com o corpo, além da apontada no numero antecedente, são duas: --- em cada corpo humano existe uma \emph{unica} alma, que é a \emph{intellectual}; --- a alma está \emph{toda} em \emph{todo} o corpo e \emph{toda} em \emph{cada} parte d'elle. --- Por isso nos dois artigos seguintes trataremos da \emph{unidade} da alma humana, e da sua \emph{séde} no corpo.

%%>>> art-10

\section*{Artigo II.\\ Unidade da alma humana.}
\addcontentsline{toc}{section}{Artigo II. --- Unidade da alma humana.}

%%>>> t49-2

\tese{49}{Estado da questão.}
O homem é dotado de vida vegetativa, sensitiva e intellectual, Logo deve existir nelle o principio d'esta triplice vida. --- Agora pergunta-se : Esse principio será triplice, ou um só? Por outras palavras: Haverá no homem tres almas realmente distinctas entre si, uma \emph{vegetativa} como a dos vegetaes, outra \emph{sensitiva} como a dos animaes, e uma terceira \emph{intellectual} propria da especie humana, --- ou haverá uma só alma, de que deriva a triplice vida? E' o que vamos investiear no presente artigo \footnote{Os philosophos, tanto antigos como modernos, occuparam-se da questão ácerca da unidade ou multiplicidade do principio vital do homem. O resultado d'estas investigações não foi uniforme. --- Uns, como \emph{Platão}, \emph{Philüo}, \emph{Plutarcho}, \emph{Galeno}, \emph{Lamarse}, admittiram no homem tres almas; --- outros, como os \emph{Manicheus}, os \emph{Apollinaristas}, \emph{Occam}, \emph{Bacon}, \emph{Gassendi}, \emph{Buffon}, \emph{Ahrens}, \emph{Jouffroy}, \emph{Maine de Biran}, \emph{Ubaghs}, \emph{Ginther}, \emph{Baltzer}, attribuem ao homem \emph{duas} almas, das quaes uma é principio da vida \emph{intellectual}, e á outra o é da vida \emph{vegetativa} e \emph{sensitiva}; --- outros, como \emph{Aristoteles}, \emph{Santo Agostinho}, \emph{Santo Thomaz}, seguidos por muitos sabios, sustentam que no homem existe uma só alma, a \emph{intellectual}, que é o principio da triplice vida.\par Presentemente, os philosophos e os naturalistas dividem-se em duas opiniões. Uns admittem no homem um principio \emph{vital}, distincto da alma \emph{intellectual}, \emph{e} o seu systema chama-se \emph{vitalismo} ou \emph{didynamismo}. Outros sustantam que a alma \emph{intellectual} é o \emph{unico} principio vital no homem, eo seu systema chama-se \emph{animismo} ou \emph{monodynamismo}. A questão, portanto, está circumscripta entre estes dois systemas. --- Alguns admittem duas almas: uma principio da vida \emph{vegetativa}, outra principio da vida \emph{sensitiva} e da \emph{intellectual}. Mas estes, que admittem no homem \emph{duas} almas, porque é que\pgn{76}não admittem \emph{tres}? Não ha, porventura, entre a sensação e a intelligencia a immensa distancia que ha entre a sensação e a nutrição? E, admittindo elles duas almas, porque é que não attribuem á alma \emph{vegetativa} a sensação, que attribuem á alma \emph{intellectual} ? --- Os philosophos anteriores a \emph{Descartes} reconheceram duas almas, uma \emph{sensitiva} e outra \emph{intellectual}. Veio \emph{Descartes}, e, admittindo apenas o \emph{corpo} (que, segundo elle, é só \emph{extensão}) e o \emph{espirito} (que, sempre na sua opinião, é só \emph{pensamento}), excluiu toda e qualquer actividade intermedia entre o pensamento e o movimento passivo, identificou a \emph{sensibilidade} com o \emph{pensamento puro}, e assim concluiu que no homem ha só a alma \emph{intellectual}, principio da vida \emph{intellectual} e \emph{sensitiva}, e que o principio da vida \emph{vegetativa} não é a alma, mas é a propria disposição dos orgãos (\emph{Cosmologia}, sec. I, c. II, a. II). Mas a theoria de \emph{Descartes} não era fundada. A observação dos phenomenos physico-chimicos demonstra que a materia não só é \emph{extensa}, mas tambem, e sobretudo, é \emph{activa}; e \& observação dos phenomenos da consciencia demonstra que a \emph{essencia} da alma não é o \emph{pensamento}, mas é essa energia profunda, essa actividade substancial, de que o pensamento e a vontade são apenas modalidades e resultancias. D' este modo, o corpo e o espirito, --- que \emph{Descartes} tinha representado como dois mundos separados por distancias insuperaveis e só relacionados por uma certa /\emph{iarmonía preestabelecida}, --- acharam-se naturalmente reunidos por uma similhanca inesperada --- a \emph{actividade corporea} e a \emph{actividade espiritual}. Mas d' aqui outra divergencia. Alguns, os \emph{vitalistas}, attribuiram ao espirito o pensamento, e ao corpo a extensão e a vida, mas sustentaram a separação radical e completa entre a vida e o pensamento, e a absoluta incompatibilidade d'estas duas operações na mesma substancia, e por isso concluiram que no homem ha duas almas: uma principio da vida \emph{vegetativa} e \emph{sensitiva} (a alma \emph{vivente}), outra principio da vida \emph{intellectual} (a alma \emph{pensante}). Outros, porém, os \emph{anünistas}, reconhecendo que as duas actividades, --- a corporea e a espiritual, --- embora essencialmente distinctas e diversas, são sempre \emph{actividades}, isto é, são sempre formas \emph{simples} e \emph{inextensas}, e por isso compativeis entre si e capazes de coexistir na mesma substancia, consideraram a vida vegetativa, sensitiva e intellectual 1o homem como tres distinctos e diversos graus de energia, provenientes da essencia de uma unica forma substancial, que é a \emph{alma intellectual}.\par Coherentes com os nossos principios, n'este artigo combatemos o \emph{vitalismo}, como os mais systemas que admittem no homem varios principios vitaes, e defendemos o \emph{animismo}. --- S. Thomaz trata d'esta questão --- na \emph{Summ}. \emph{theol}., I p., q. (6, a. 8; na \emph{Sum}. \emph{contra gent.,l}. IL, c. 58; --- nas \emph{Quest}. \emph{Disp}. \emph{De anima}, aa. 9 e 11, \emph{De spirit}. \emph{creat}. a. 3.}.

%%>>> t50

\pg{76} \tese{50}{O homem tem uma só alma.}

a) O homem a dotado de unidade \emph{substancial}; pois que a alma e o corpo, pela sua união, constituem uma só substancia completa, Esta unidade deriva ou da alma ou do corpo; visto serem estes os elementos essenciaes do homem. Mas não deriva do corpo; porque a materia é por si multipla, divisivel e tende a desaggregar-se. Logo deriva da alma. Se deriva da alma, esta deve ser uma só; porque, se fossem duas ou mais, \pg{77}  nãd teriamos no homem um só ser vivo, mas dois ou mais seres vivos, o que é absurdo. Logo o homem tem uma só alma \footnote{Este argumento, baseado na \emph{unidade} do ser humano, é convincente. O homem não é o \emph{aggregado} de tres substancias completas: uma vegetativa, outra sensitiva, e outra intellectual; mas é uma só substancia, ainda que composta, dotada de triplice grau de vida. Ora, se no homem existissem tres principios formaes, distinctos entre si, o homem seria um \emph{aggregado} de tres substancias \emph{completas}, e a sua unidade seria apenas \emph{accidental}, e nao \emph{substancial}. Por tanto, se o homem é um só ser, uma só substancia completa, um sô tambem deve ser o principio formal, pelo qual elle vegete, sinta e entenda. Diz S. Thomaz: « Do principio, de que procede o ser, deve proceder a unidade, porque a unidade é uma consequencia do ser, e é o proprio ser. Por isso a forma, que dá o ser, deve dar tambem a unidade. Logo se no homem existissem duas ou mais almas, como formas distinctas, o homem não seria um só ente, mas um aggregado de entes: Ab \emph{eodem aliquid habet esse et unitatem}; \emph{unwm enim consequitur ad ens}. \emph{Cum igitur} a \emph{forma unaquaeque ves habeat esse}, \emph{a}. \emph{forma etiam habebit}. \emph{unitatem}. \emph{Si igitur in homine ponuntur plures animae}, \emph{sicut}. \emph{diversae formae}, \emph{homo non erit unum}. \emph{ens}, \emph{sed}. \emph{plura} » (\emph{Sum}. \emph{contr}. \emph{gent}., l. Yl, c. 58). --- Se, portanto, distinguimos no homem a vida vegetativa da vida intellectual, esta distineção não suppõe substancias realmente diversas, porque não podemos separar, no homem, a substancia pensante da substancia vivente, mas indica apenas diversas energias da mesma forma, que é o principio primeiro e radical de toda a actividade do ente, isto é, suppõe apenas uma diversidade nos principios secundarios e proximos de operações.\par Podem objectar: Embora se admittissem no homem varios principios vitaes, comtudo haveria sempre uma certa unidade, porque poderiam esses principios estar por tal modo coordenados entre si, que um fosse perfeição e complemento do outro. --- Resposta. Essa unidade seria apenas de \emph{ordem} (que é a ultima entre as especies de unidade) e por isso seria sempre \emph{accidental}; ao passo que á unidade entre as tres vidas é muito \emph{intima}, é verdadeiramente \emph{substancial}, porque se refere sempre ao mesmo sujeito, ao mesmo ser vivo, sendo sempre um e o mesmo homem que vive, sente e entende (Cf. \emph{Contra Gentes}, l. c.).}.

b) Se o homem tivesse tres almas distinctas, e por isso tres principios de operações vitaes, estas operações não poderiam impedir-se mutuamente; porque não são contrarias, mas apenas diversas entre si. Ora as operações da triplice vida humana impedem-se mutuamente, de modo que, se uma é intensa, outra é remissa. Logo o homem tem uma só alma \footnote{\emph{A} experiencia mostra que uma applicação constante e intensa das faculdades intellectuaes é nociva ás faculdades nutritivas e pode impedir o regular exercicio dos sentidos; como tambem o immoderado uso dos prazeres sensiveis torna imbecil a intelligencia e fraca a vontade. Ora estes phenomenos denotam que um só é o principio radical, de que brota toda a energia humana, de modo que esta energia não pode superabundar numa faculdade, sem que, ao mesmo tempo, diminua noutra. --- \emph{A}! mesma conclusão\pgn{78} chegamos, considerando a influencia que as faculdades superiores do homem exercem nas inferiores, e vice-versa; assim um intenso affecto da vontade desperta uma paixão no appetite sensitivo, como uma paixão vehemente do appetite sensitivo arrasta a vontade para o mal, que esta tinha antes detestado. Ora estes phenomenos seriam inexplicaveis, se no homem não fosse um e o mesmo o principio primeiro e radical de todas as faculdades. Diz S. Thomaz: « Cum operatio unius potentiae fuerit intensa, impeditur alterius operatio, et e contra fit redundantia ab una potentia in aliam, quod non esset, nisi omnes potentiae in una essentia animae radivarentur » (\emph{De anima}, a. 1. --- \emph{Of}. \emph{Summ}. \emph{theol}., LI p., q. 16, a. 3).}.

c) O principio da triplice vida do homem não póde ser triplice, mas deve ser um só. Nu verdade, a consciencia attesta que em nós o sujeito que entende é o mesmo que sente e vegeta. Ora, se no homem existissem tres principios de vida, o sujeito que entende, por ser distincto do sujeito que sente e do que vegeta, não poderia nem sentir nem vegetar; visto que os actos vitaes são \emph{immanentes}, isto é, devem ter o termo no mesmo sujeito em que têm o começo. Logo o homem tem um só principio da triplice vida. O principio da vida diz-se \emph{aZma}. Logo o homem tem uma só alma \footnote{E' á consciencia que nos diz que a unica nossa alma é o principio da triplice vida. --- Antes de tudo, a consciencia attesta a identidade do principio radical, de que derivam os actos da vida \emph{intellectual} e os da vida \emph{sensitiva}. E' sempre o mesmo \emph{eu} que pensa, quer, sente e soffre, e, embora o eu não exerça \emph{do mesmo modo} as operações \emph{intellectuaes} e as \emph{sensitivas}, comtudo a correlação do pensamento com a sensibilidade é tão intima o tão estreita, que estas duas series de phenomenos costumam andar sempre unidas. Uma simples reflexão sobra nós mesmos convence-nos de que o pensamento anda sempre acompanhado de imagens, e que as imagens e as ideas se excitam e provocam mutuamente. --- Diz-nos tambem a consciencia que do mesmo principio radical, de que derivam as funcções \emph{intellectuaes} e \emph{sensitivas}, derivam tambem as \emph{vegetativas}. De facto, muitas funcções \emph{vegetativas}, necessarias á conservação, nutrição e reproduccüo do individuo, são tão \emph{voluntarias}, como são \emph{voluntarias} as operações \emph{sensitivas}. Ora toda esta surprehendente correlação de todos os phenomenos da vida humana, que a consciencia nos manifesta, ó uma prova convincente da unidade e identidade do principio radical, quo é a alma.\par Esta convicção dos homens relativa á unidade e identidade da alma não é uma illusão, mas é o resultado de uma reflectida analyse dos nossos actos. E se, ás vezes, uma pessoa não tem consciencia d'essa unidade e identidade, --- ou se, por causa de um certo estado pathologico, parece desdobrar a sua personalidade e a sua consciencia, nem por isso essa unidade e identidade deixa de existir. Uma coisa é o \emph{eu}, outra coisa é a \emph{idéa} do \emph{eu}. O eu é fixo e immutavel; a \emph{idéa} do \emph{eu} pode variar, é mutavel, como variaveis são os estados psychicos. Ora a verdade das coisas deve \&preciar-se, não pelo que parece ás pessoas inconscientes ou allucinadas, mas pelo que parece ás que têm juizo. Por isso, os que dizem que a pathologia descobriu\pgn{79}no homem muitas personalidades, muitas consciencias, confundem o \emph{eu sujeito} (que, como dissemos, é indivisivel e permanente) com os \emph{estados} ou com as \emph{apreciações} do \emph{eu}, e a \emph{personalidade} com a consciencia \emph{actual} d" ella.\par Examinemos agora as principaes objecções contra esta these.\par a) Dizem: O homem define-se \emph{animal racional}. Ora o genero \emph{animal} desume-se da alma \emph{sensitiva}, e a differença \emph{racional} desume-se da alma \emph{intellectual}. Logo a alma \emph{sensitiva} no homem differe da alma \emph{intellectual}.\par Resposta. Na definição do homem o genero \emph{animal} não se desume da alma \emph{sensitiva}, nem a differença \emph{racional} da alma \emph{intellectual}, como se á alma \emph{sensitiva} e a \emph{intellectual} fossem realmente distinctas; mas o genero \emph{animal} desume-se da alma \emph{sensitiva}, isto é, da \emph{materia}, emquanto esta recebe da alma \emph{intellectual} o grau de vida \emph{sensitiva}, e a differenca \emph{racional} desume-se da mesma alma \emph{intellectual}, emquanto esta é o principio de um grau mais excellente de perfeição, isto é, da vida intellectual, e por isso não se segue que se deva admittir tal distincção. --- Como se vê, á difficuldade confunde a ordem \emph{logica} com a \emph{real}. Num e mesmo ente \emph{ontologico}, ou \emph{real}, a nossa intelligencia pode considerar muitos e diversos graus de perfeição. Mas, do facto de o homem distinguir num ente muitos e diversos graus de perfeição não se pode concluir que a distinecüo seja tambem \emph{real}, isto é, que exista nas coisas realmente e independentemente da consideração mental. Assim, contendo \emph{virtualmente} a alma intellectual tudo o que tem a alma sensitiva, pode considerar-se separadamente o grau de vida sensitiva, como a parte \emph{imperfeita} e \emph{material}, e, sendo este grau commum ao homem e ao animal, pode com elle formar-se o \emph{genero}; e o grau de vida intellectual, no que a alma \emph{intellectual} excede a \emph{sensitiva}, pode considerarse como a parte \emph{formal} e \emph{complementar}. Mas a distineção entre a alma intellectual considerada no seu mais excellente grau de vida e a mesma alma considerada no grau de vida \emph{sensitiva} é uma distincção \emph{logica}, nao se encontra realmente na alma, ainda que nº ella tenha o seu fundamento. (S. Thom., \emph{Sum}, \emph{Theol}., \emph{I} p., q. 76, a. 8 ad 4) --- Por outras palavras: A alma emquanto é \emph{espiritual} e excede a capacidade da materia, distingue-se de si mesma, emquanto é acto ou forma do corpo; mas essa distinecüo é apenas \emph{logica}, é o resultado da consideração da intelligencia.\par Pp) Continuam: Um ente não póde estar em lucta comsigo mesmo. Ora, no homem, o appetite sensitivo está em lucta com a vontade. Logo estas duas faculdades não derivam do mesmo principio. :\par Resposta. Concedemos que um ente não póde estar em lucta comsigo mesmo, quando não possue faculdades de ordem diversa; mas negamos que não possa estar em lucta, quando é dotado de faculdades de ordem diversa, porque, então, o bem de uma faculdade póde ser contrario ao da outra;\pgn{80}d' ahi a lucta. E' o que acontece no homem. Elle é dotado de faculdades de ordem diversa; o appetite \emph{sensitivo} pertence á ordem \emph{sensitiva}, e a \emph{vontade} partence á ordem \emph{Intellectual}. Pertencendo a ordens diversas, essas duas faculdades têm por termo objectos diversos e ás vezes contrarios, de modo que o que é agradavel ou util a uma, é desagradevel ou nocivo a outra. --- Accrescente-se que esta lucta, longe de provar a \emph{multiplicidade} do principio vital, prova, pelo contrario, a sua \emph{unidade}, Com effeito, cada um de nós experimenta que a lucta entre o appetite sensitivo e a vontade se trava \emph{num e mesmo sujeito}, que é agitado por inclinações e tendencias contrarias; o que não aconteceria, se aquellas faculdades não derivassem do mesmo principio, mas sim de dois principios distinctos, pois então cada um operaria dentro da sua esphera, sem se impedirem mutuamente e luctarem entre si.\par c) Proseguem: Oparações \emph{especificamente} diversas suppõem principios \emph{especificamente} distinctos. Ora, no homem, existem operações \emph{especificamente} diversas, --- as vegetativas, as sensitivas e as intellectuaes. Logo existem nelle tres principios \emph{especificamente} distinctos.\par Resposta. A diversidade \emph{especifica} das operações suppõe uma diversidade \emph{especifica} no principio, quando se trata de individuos numericamente distinctos. Mas, quando se trata de um e o mesmo individuo, aquella diversidade de operações só denota uma diversidade de faculdades no mesmo principio, que, sendo unico na essencia, é multiplo na força ou energia.\par d) Insistem: Se, no cadaver do animal, se appliear a centelha electrica ás extremidades dos nervos, suscitam-se movimentos de \emph{contractilidade}; e similhantes movimentos continuam por algum tempo numa cabeça decepada, ou nos membros que se amputaram. Ora estes movimentos são \emph{vitaes}, e todavia a alma ja não reside no cadaver, nem nos membros cortados. Logo é necessario admittir no homem um principio \emph{vital}! distincto da alma.\par Concedemos os factos, mas negamos a sua apreciação e a consequencia que d' elles tiram os adversarios. Na verdade, emquanto aos movimentos de \emph{contractilidade}, suscitados pela centelha electrica, devemos dizer que elles são \emph{mechanicos} e não \emph{vitaes}, porque derivam de um agente \emph{extrinseco}; e por isso não se segue que deva admittir-se no homem um principio vital distincto da alma. Emquanto aos movimentos de \emph{contractilidade}, que por algum tempo continuam nas cabeças dos guilhotinados, ou nos membros recentemente amputados, dizemos que elles derivam da acção exercida pela alma nos musculos do homem antes da sua morte, e não cessam immediatamente depois da morte, assim como não cessam immediatamente as vibrações das cordas de uma harpa, quando os dedos deixaram de a tocar. rj pois, falso que exista no homem um principio \emph{vital} distincto da alma.}.

%%>>> t51

\pg{78} \pg{79} \tese{51}{A unica alma do homem é a intellectual.}

a) A unica alma do homem é a forma \emph{substancial}, que o colloca na especie dos seres corporeos \emph{intelligentes}, distinguindo-o dos seres de outras especies. Ora a forma \emph{substancial}, que colloca o homem na especie dos seres corporeos \emph{intelligentes}, é a alma \emph{intellectual} ; pois cada effeito exige uma causa proporcionada. Logo a unica alma do homem é a \emph{intellectual}.

b) O homem tem uma só alma, e é dotado de vida vegetativa, sensitiva e intellectual. Logo a unica alma, existente no homem, 6, ao mesmo tempo, principio da sua triplice vida. Ora a alma, que é o principio da triplice vida, ou é a \emph{vegetativa}, ou a \emph{sensitiva}, ou a \emph{intellectual} ; pois ha só estas especies de alma. Mas não é a \emph{vegetativa}; porque esta não póde ser principio da vida \emph{sensitiva} e muito menos da \emph{intellectual}, aliás o vegetal deveria possuir sensibilidade e entendimento. Nem \pg{81}  é a \emph{sensitiva}; porque esta não póde ser principio da vida \emph{intellectual}, aliás todo o animal seria dotado de intelligencia. Logo a unica alma do homem é a \emph{intellectual} \footnote{Objectam: A alma humana é simples e indivisivel, e as funcções vitaes são muitas e diversas. Logo a alma não póde ser o principio de todas as funcções vitaes do homem.\par Respondemos que á variedade das funcções vitaes póde e deve conciliar-se com a simplicidade da alma; por quanto á alma, embora seja simples na essencia, todavia é multipla nas forças ou faculdades, como de um só centro indivisivel partem muitos e diversos raios. --- D'hai se vê o erro d'alguns escriptores, que, pelo facto de as nossas faculdades serem muitas e diversas, concluem que no homem o principio vital é distincto da alma.}.

%%>>> t52

\pg{80} \tese{52}{A alma intellectual a principio tambem da vida vegetativa e sensitiva no homem.}

a) As formas \emph{substanciaes} estão por tal modo subordinadas entre si, que á forma superior possue toda a virtude das inferiores, e por isso aquella communiea ao corpo a perfeição, que estas lhe teriam communicado. Ora a alma \emph{intellectual} é forma \emph{substancial} do corpo e é superior á alma \emph{vegetativa} e á \emph{sensitiva}. Logo a alma \emph{intellectual} communica ao corpo toda a perfeição, que a alma \emph{vegetativa} e a \emph{sensitiva} lhe teriam communicado. Mas a alma \emph{vegetativa} teria communicado ao corpo a vida \emph{vegetativa}, e a alma \emph{sensitiva} teria communicado a vida \emph{sensitiva}. Logo a alma \emph{intellectual} communica ao corpo a vida \emph{vegetativa} e a \emph{sensitiva}, e, por isso, ella é principio tambem da vida \emph{vegetativa} e \emph{sensitiva} no homem.

b) A forma \emph{substancial} é O principio de \emph{toda} a actividade, que uma substancia composta possue. Ora a alma \emph{intellectual} é a forma \emph{substancial} do homem. Logo a alma \emph{intellectual} é O principio de \emph{toda} a actividade, de que o homem é dotado. Ora a actividade do homem não só abrange a vida intellectual, mas tambem a vegetativa e a sensitiva. Logo a alma \emph{intellectual} é principio tambem da vida \emph{vegetativa} e \emph{sensitiva} no homem \footnote{A razão, pela qual a alma humana não só é principio das faculdades \textbar{} intellectuaes, mas tambem das sensitivas e vegetativas, encontra-se na ex- \textbar{} cellencia da propria alma, a qual, por ser a mais nobre de todas as formas substanciaes, contém, na sua unica essencia, todos os graus inferiores de entidade. Nem isto é contrario á simplicidade da alma; pois aquella multiplicidade de graus de entidade não denota na alma uma extensão \emph{quantitativa} de partes materiaes, mas uma extensão \emph{qualitativa} de força ou actividade. Por isso a alma intellectual exerce no homem as mesmas funcções, \textbar{} que a alma vegetativa exerce nas plantas e a sensitiva nos animaes, e, além d'isso, é principio da vida intellectual. :\par Ougamos o Doutor Angelico: « A mesma alma no homem é intellec- \textbar{}\par 6 Vor. IL \textbar{}\par \textbar{} \textbar{}\pgn{82}tual, sensitiva e nutritiva. Isto percebe-se facilmente, se se considere a differenga das especies e das formas. Na verdade, as especies, e por isso os seus principios formaes, differem uma de outra, conforme o maior ou menor grau da sua perfeição; assim, no mundo, os seres vivos são mais perfeitos que os seres não-vivos, e os animaes mais que as plantas, e os homens mais que os animaes; porque o grau de perfeição, que um d' estes seres encerra, differe do grau de perfeição, que encerra o outro. Por isso Aristoteles (\emph{De anima}, lib. II, text. 30 et 81) compára as especies das coisas aos numeros, que differem especificamente entre si pela addição ou subtraecüo de uma unidade, e compára tambem as diversas almas ás diversas especies de figuras, das quaes uma contém a outra, como o pentagono contém o tetragono, e o excede. Do mesmo modo, a alma intellectual contém \emph{virtualmente} tudo o que possuem a alma sensitiva dos animaes e a alma nutritiva das plantas. Por isso, assim como uma superficie de figura pentagona não é tetragona por uma figura e pentagona por outra, porque seria superflua a figura tetragona por estar contida na pentagona, assim tambem Socrates não é homem por uma alma e animal por outra, mas é homem e animal por uma e mesma alma » (\emph{Summ}. \emph{theol}., p. I, q. 16, a. 3.\par O Dr. \emph{Frédault} reconhece tambem que uma multiplicidade de almas ou de principios vitaes é, pelo menos, inutil, e só serviria para tornar inexplicavel a harmonia das partes e a unidade das operações, que se admiram no homem. Eis as suas palavras: « Se diversos principios coexistissem juntamente no homem, seria necessario --- que um dirigisse os outros, para assim haver unidade, --- que os regulasse em conformidade com as suas formas, com o seu grau, com a sua ordem, --- que nenhum d'elles podesse operar independentemente do outro. Ora é claro que, para haver esta subordinação, seria preciso quo o principio superior estivesse, a0 mesmo tempo, no todo e em cada uma das partes. E se este principio esta em toda a parte, para mover tudo, dirigir tudo, unificar tudo, que necessidade temos nós de admittir outros principios? E se, apezar de tudo, quizerem admittilos, quem não vé que elles, só por si, são nada, e que são completamente aniquilados pelo principio dominante, visto que, sem este principio, não podem ter acção, nem direcção, nem existencia util? » (\emph{Ant/mopologie}, pag. 195).\par Esta theoria da subordinação das formas concilia admiravelmenite a variedade com a unidade, a indivisibilidade do composto humano com a prodigiosa diversidade das suas potencias.\par Examinemos algumas difficuldades contra a these.\par a) Dizem: Se a alma intellectual fosse tambem principio das operações vegetativas, deveria ter consciencia d'ellas. Ora não a tem. Logo.\par Resposta. Antes de tudo, não é verdade que a alma não tenha consciencia de nenhuma operação vegetativa; porque os actos da nutrição, da respiração, da reproducção, que constituem uma grande parte da vida organica, não só são percebidos pela consciencia, mas até estão sujeitos ao dominio da vontade. E se uma outra serie de actos da vida organica escapa á percepção\pgn{83}da consciencia, como a digestão, a chylificacao, a circulação do sangue, a secreção da bilis, isto não quer dizer que esses actos não procedam, como de principio \emph{radical}, da alma intellectual (aliás deveriamos admittir só para \& vida vegetativa duas almas), mas quer dizer que esses actos, embora derivem \emph{radicalmente} da alma, todavia \emph{prowimamente} derivam do \emph{composto} e por isso, não estando ligados, de modo algum, com as operações intellectuaes, escapam á percepção da consciencia. Só o composto pode perceber as operações do composto. Deus não dotou o homem de um orgão central, que recolhesse esses actos da vida organica, como o dotou de um orgão central para a percepção dos actos dos sentidos particulares, porque a consciencia de certos actos da vida organica, além de inutil, seria prejudicial e origem de innumeraveis perplexidades, e tornar-se-hia muito diffieil, senão impossivel, o exercicio da nobilissima vida intellectual.\par b) Acerescentam: Se a alma intellectual não se distinguisse da alma sensitiva nem da vegetativa, seguir-se-hia que a materia, assim como recebe da alma a vida vegetativa e a sensitiva, assim tambem deveria receber a vida intellectual. Ora o consequente é falso; aliás o acto do entendimento seria organico, e não espiritual. Logo a alma intellectual distingue-se da alma vegetativa e da sensitiva, e por isso não é nem póde ser principio da vida vegetativa e da sensitiva no homem.\par Resposta. A alma humana, come dissemos, por ser a mais elevada forma \emph{substancial}, contém \emph{virtualmente} na sua essencia todos os graus de entidade das formas inferiores, e por isso da alma vegetativa e da sensitiva. Quando a alma se une ao corpo, só lhe communica aquelles graus de entidade, que o corpo pode receber. Ora o corpo, se póde receber os graus de vida vegetativa e. sensitiva, pois o sujeito que vegeta e sente é o composto de materia e de forma, não póde receber o grau de vida intellectual, porque a vida intellectual excede, pela sua excellencia e universalidade, as condições da materia. Por isso a alma humana, na sua união com o corpo, communica aos orgãos as faculdades da vida vegetativa e sensitiva, mas reserva para si as faculdades da vida intellectual. (Cf. c. \emph{Gentes}. l. II, c. 69).\par 0) Continuam: As forças \emph{vegetativas} e \emph{sensitivas} são \emph{materiaes}. Ora forças \emph{materiaes} não podem derivar de um principio \emph{espiritual}, como é a nossa alma.\par Respondemos que as forças \emph{vegetativas} e \emph{sensitivas} são \emph{materiaes} emquanto residem num \emph{sujeito material}, mas não são \emph{materiaes} emquanto são \emph{materia} ou compostas de materia, pois são \emph{simples}. Sendo \emph{simples}, \textbf{essas} forças podem derivar da nossa alma, que é essencialmente \emph{simples}.\par De resto, a doutrina, que sustentamos, é defendida por muitos distinctos naturalistas, como Sthal (\emph{Theoria medica nova}), Brillier (\emph{Du principe vital et de Váme pensante}), Frédault (\emph{Traité de anthropologie}). Limitamo-nos a citar as palavras de Chauffard (\emph{La} Vie): « A alma é a \emph{universal} causa \emph{humana}, 0 \emph{principio de todas} as \emph{nossas actividades} n.}.

%%>>> t53

\pg{82} \tese{53}{A alma intellectual é a unica forma substancial no homem.}

a) A forma \emph{substancial} é a realidade que dá a \emph{unidade} ao ente e o constitue numa determinada especie. Ora, se no homem existisse, além da alma intellectual, uma outra forma \emph{substancial}, \pg{83}  o homem não seria dotado de unidade e pertenceria, ao mesmo tempo, á duas ou mais especies, o que é absurdo. Logo a alma \emph{intellectual} é a unica forma \emph{substancia}] no homem.

b) Se, além da alma \emph{intellectual}, existisse no homem uma outra forma \emph{substancial}, que o constituisse na especie de \emph{ente corporeo}, --- não só a alma communicaria um ser \emph{accidental}, \pg{84}  á; poisque toda a perfeição, que se aecrescenta a um ser anteriormente constituido, é \emph{accidental}, --- mas tambem a sua união

\emph{z} com o corpo não seria \emph{substancial}, mas apenas \emph{accidental}, porque a união desume do termo a propria especie, e neste caso o termo não seria uma \emph{substancia} mas um \emph{aggregado} de substancias, e a união no \emph{aggregado é accidental}. Logo a alma \emph{intellectual} é á unica forma \emph{substancial} existente no homem \footnote{Nos numeros anteriores dissemos que a alma intellectual, contendo em si a virtude, a perfeição da alma vegetativa e da alma sensitiva, communica ao homem os graus da vida vegetativa, os da sensitiva e mais alguma coisa, isto é, a vida intellectual. --- Restava-nos ainda ver se a alma intellectual, além de excluir do composto humano toda e qualquer outra alma, exclua tambem toda e qualquer outra forma substancial, de modo que ella communique ao composto não só o que lhe poderiam communicar as outras almas, mas tambem o que lhe poderiam communicar as outras formas. Por outras palavras: restava-nos investigar se a alma intellectual se une à materia \emph{prima}, ou se une ao corpo já constituido no ser de corpo por outra forma substancial, a qual, juntamente com a alma, permaneça actualmente no homem. --- S. Thomaz, seguindo a doutrina de Aristoteles, ensina que em todo e qualquer composto substancial existe uma só forma substancial, unida por si e immediatamente á \emph{materia prima}, e que por isso o homem, pela alma intellectual, não só é homem, mas é animal, é vivo, é corpo, é substancia, é ente: « In homine non est alia forma substantialis, quam anim» intellectiva, et per eam homo non solum est homo, sed et animal et vivum et substantia et ens » (\emph{De Spirit}. \emph{creat}., a. 3). --- O Doutor Angelico trata d'esta questão tambem --- na \emph{Sum}. \emph{theol}, p.l, q. 76, a. 4, --- \emph{De anima}, a. 9; --- \emph{c}. \emph{Gentes}, l. II, c. 58, e no opusc. \emph{De plural}. \emph{formarwum}.\par A razão d'esta importante verdade é. sempre a mesma: a \emph{equivalencia virtual} das formas: a forma superior contém \emph{virtualmente} a inferior, e por isso pode substituil-á com vantagem. Diz S. Thomaz: « Forma perfectior \emph{virtute} continet quidquid est inferiorum formarum; et ideo una et eadem existens perficit materiam secundum diversos perfectionis gradus. Una enim et eadem forma est per essentiam, per quam homo est ens actu, et per quam est corpus, et per quam est vivum, et per quam est animal, et por quam est homo » (\emph{Sum}. \emph{Theol}. p. I, q. 76, a. 6 ad 1). Por isso, a nossa alma, embora, \emph{formalmente} e pela sua essencia, seja incorporea e inextensa, todavia, \emph{virtualmente}, emquanto contém a virtude ou perfeição das formas inferiores, é corporea e extensa, e, communicando o proprio ser á materia, torna o composto \emph{formalmente} corporeo e extenso. --- Se uma forma contém a energia da inferior, é claro que uma dispensa a outra. E' \emph{o} methodo da \emph{economia}, que a Escholastica exprimiu com estas palavras: não se devem multiplicar as causas sem necessidade --- \emph{entia non sunt multiplicanda sine necessitate}. Este axioma é admittido como indiseutivel pelos mais sabios naturalistas. Newton diz: « Não se devem admittir mais causas que as que são necessarias para explicar os phenomenos. A natureza não faz nada em vão,\pgn{85}e seria fazer coisa inutil realizar por um grande numero de causas o que se poderia conseguir por um numero mais pequeno » (\emph{Principios mathematicos de philosophia natural}, l. 3).\par Dissemos que a alma dá o primeiro ser, o ser substancial ao corpo, isto é, faz com que o corpo seja \emph{corpo}. Esclarecamos melhor este ponto, ainda que se deva fazer alguma repetição. A \emph{materia} é uma verdadeira realidade, mas potencial, incompleta, indeterminada, e por isso não pode existir só “por si; e se existe, precisa de um principio formal que lhe communique o acto e a especificação, e a complete e determine. O composto d'estes dois elementos é o \emph{corpo}. Portanto a materia, só por si, não é propriamente \emph{corpo}; porque a propria extensão, de que ella é raiz, não pode conceber-se sem unidade, nem tornar-se sensivel sem acção; e a unidade e a acção derivam de um principio simples. Ora a alma, unindo-se á materia, exerce o officio d'esse principio simples, e por isso communica-Ihe o primeiro ser substancial de \emph{corpo}. --- D'onde se vê --- 1) que a alma, embora não seja \emph{formalmente} corporea, todavia é tal \emph{virtualmente}; porque, embora não seja corpo, todavia a alma faz com que seja corpo a materia, a que ella se une como forma substancial (Cf. S. Thom., \emph{De spirt}. \emph{creat}., a. 3j; --- \emph{De anima}, a. 11); --- 2) que, propriamente, é a alma que contém o corpo, e não viceversa: \emph{magis anima continet corpus}, \emph{quam e converso} (\emph{Sum}. \emph{theol}, p. I, q. 16, a. 8); porque é ella, que mantém todas as partes do corpo na unidade do ser e da vida; --- 8) que, quando se diz que a alma é \emph{acto} ou \emph{forma} do corpo, devemos entender que o corpo é corpo pelo facto mesmo de a alma Ihe se unir como acto ou forma, assim como quando dizemos que a luz é forma do lucido, entendemos que o sujeito é lucido só pelo facto de receber em si a luz (\emph{Sum}. \emph{Theol}. p. I, q. 76, a. 4 ad 1): --- 4) que, quando dizemos que \emph{o homem é composto de alma e de corpo}, não devemos entender que o corpo seja corpo independentemente da alma, mas --- ou que \emph{corpo} significa a \emph{materia}, --- ou que o \emph{corpo}, isto é, a materia informada pela alma e por esta determinada ao ser substancial de corpo, se considera como \emph{sujeito} e \emph{potencia} em relação á mesma alma, em quanto a alma é principio de vida (\emph{De spirit}. \emph{creat}., a. 8, ad 2).\par \emph{Scoto} admittiu nos seres vivos, além da alma, uma forma de. \emph{corporeidade}, que dispõe a materia para receber a alma, e que por si não constitue o corpo num determinado genero ou especie. --- Esta opinião é hoje seguida pelos \emph{atomistas} e \emph{dynamistas}, os quaes afirmam que o nosso corpo, independentemente da alma, é uma substancia completa e perfeita.\par A doutrina de \emph{Scoto} é falsa; porque, se a materia tivesse a forma de \emph{corporeidade}, distincta da alma, a alma não seria forma \emph{substancial} mas \emph{accidental}. --- Além d'isso, é impossivel que exista um \emph{corpo}, que não pertença a uma certa e determinada \emph{especie}.\par Objectam: Se a alma faz com que o nosso corpo seja \emph{corpo}, este, independentemente da alma, é nada.\par A objecção confunde a \emph{materia} com o \emph{corpo}, quando é certo que a materia é \emph{um} dos elementos \emph{substanciaes} do corpo, pois este é composto de\pgn{86}\emph{materia} e de \emph{forma}. A \emph{materia} não é \emph{actualmente} corpo, senão emquanto está unida á \emph{forma}. Por isso, é falso que o nosso corpo, independentemente da alma, seja \emph{nada}; pois que é \emph{materia}, isto é, é uma realidade positiva. --- A \emph{materia}, que por si é incompleta e determinavel, torna-se completa e determinada pela alma, e por esta e com esta constitue o \emph{corpo}.\par Insistem: Se a alma exercesse tambem as funcções de forma de \emph{corporeidade}, deveria admittir-se no corpo, privado da alma, uma forma \emph{cadaverica}.\par Concedemos tudo o que os adversarios objectam. Mas d'hai derivará algum absurdo? Certamente que não. A materia não póde existir sem forma substancial Por isso, quando a alma intellectual se separa do corpo, este não fica sem forma \emph{substancial}. E, como o corpo separado da alma é e diz-se \emph{cadaver}, a forma \emph{substancial} que succede á alma é a \emph{cadaverica}; pois é a forma que determina a materia. --- A forma cadaverica conserva no corpo morto o mesmo organismo e a mesma disposição de partes, pois a forma deve ser proporcionada á materia; mas esta conservação é passa geira, porque a forma \emph{cadaverica} não é uma forma para a qual, como para seu termo, tende a natureza, mas é apenas uma \emph{transição} entre a alma e as formas dos elementos que fazem parte do composto humano.\par Pergunta-se: Qual é a causa productora da forma \emph{cadaverica}?\par A causa productora da forma \emph{cadaverica} deve procurar-se nas qualidades dos corpos elementares, que compõem o nosso organismo. Essas qualidades, depois de se separar do organismo o principio \emph{vital} que as subordinava, recuperam a sua independencia e, operando, reproduzem nos mesmos corpos elementares as formas \emph{substanciaes}, que elles possuiam antes de fazerem parte do organismo humano. Ora, assim como na geração a força, que está no germen, antes de produzir a ultima forma perfeita, produz varias formas intermedias e imperfeitas; assim tambem as qualidades dos corpos elementares, por isso mesmo que operam para resolver o corpo vivo nos seus elementos, produzem varias formas intermedias, começando pela forma, que corresponde ás disposições deixadas no organismo pelo principio \emph{vital} e que é á forma \emph{cadaverica}.}.

%%>>> t54

\tese{54}{Corollario.}
Das theorias expostas resulta que a alma intellectual não é uma e a mesma em todos os homens. Na \pg{85}  verdade, a alma é o principio do ser substancial no homem. Ora o ser substancial não é um e o mesmo em todos os homens, mas o ser substancial de um é distincto do ser substancial dos outros. Logo a alma intellectual não é uma e a mesma \pg{86}  em todos os homens. Devemos, pois, dizer que ha tantas almas, quantos são os homens, e que ellas, embora pertençam todas á mesma especie, são comtudo \emph{numericamente} distinctas \footnote{Averroes, philosopho arabe (1120-1198 ou 1206), nascido em Cordova de Hespanha, disse que cada um dos homens possue uma alma \emph{sensitiva} propria e distincta da dos outros, mas que a alma \emph{intellectual} é uma e a mesma em todos os homens. Esta unica alma intellectual, segundo \emph{Averroes}, existe fora dos homens, mas illumina-os a todos, e, segundo a diversidade das especies que existem na imaginação de cada um, produz diversos pensamentos, que d' este modo se tornam proprios de cada um; assim como o mesmo ar, penetrando nos tubos de um orgão, produz diversos sons segundo a diversidade dos mesmos tubos.\par A opinião de Averroes é \& dos \emph{pantheistas} de todos os seculos, e principalmente dos philosophos \emph{transcendentaes} da Allemanha, --- \emph{Hegel}, \emph{Schelling}, \emph{Fichte}, que --- ou admittem uma unica realidade, o \emph{eu-puro}, --- ou identificam o sujeito com o objecto. Mas essa opinião é falsa pelas seguintes razões:\par \emph{a}) \emph{Ente} e uno são, na realidade, uma e á mesma coisa. Logo o principio,\pgn{87}que communica a um ser a entidade, communica-lhe tambem a unidade. Mas o principio, que communica ao homem a entidade, constituindo-o numa determinada especie e distinguindo-o de todos os outros seres, é a alma intellectual. Logo tambem a alma intellectual é o principio que communica ao homem a unidade, pela qual elle não está dividido em si, mas está dividido de todos os outros. Ora a alma intellectual não poderia communicar ao homem a unidade, se não a possuisse; tanto mais que a alma confere a unidade pelo facto de se dar a si mesma. Logo a alma intellectual é dotada de unidade, pela qual não está dividida em si, mas está dividida de todas as outras almas. Logo a alma intellectual não é uma e a mesma em todos os homens. Em poucas palavras: A alma é a forma que dá o ser; ora o ser de um homem é distincto do ser dos outros homens; logo a alma de um homem é distincta da alma dos outros. Diz S. Thomaz: « Impossibile est plurium numero diversorum esse \emph{unam} formam, sicut impossibile est quod eorum sit unum esse; nam forma est essendi principium ». (\emph{Sum}. \emph{Theol}, p. I, q. 66, a. 2; Cf. c. \emph{Gentes}, lib. IL, e. 78).\par \emph{D}) A consciencia attesta a cada homem que as percepções intellectuaes, que produz, pertencem a elle e não são communs aos outros homens. Ora as operações derivam da alma, como do seu principio, e o principio de uma operação pertence a quem pertence a propria operação. Logo cada homem possue uma alma propria e distincta da alma dos outros homens, e por isso a alma não é uma e a mesma em todos os homens. --- Averroes responde que a nossa consciencia nos attribue as percepções intellectuaes pelo facto de nos pertencerem as imagens sensiveis, ás quaes unir-se-hia a alma intellectual. Mas esta resposta é ineffioaz. As imagens sensiveis, como veremos, encerram o \emph{objecto} da. percepção intellectual. Ora do facto de o objecto da percepção pertencer a um certo individuo não se póde concluir que ao mesmo individuo deva pertencer a percepção d aquelle objecto. Concluimos, pois, que as almas intellectuaes, embora pertencam todas á mesma especie, são comtudo \emph{numericamente} muitas e distinctas entre si (Cf. \emph{Sum}. \emph{theol}., p. I, q. (6, a. 1; --- \emph{De unitate intellectus contra Averroistas}).\par *\par Pergunta-se: Qual é o principio que torna \emph{individua} a alma humana ?\par Respondemos quo o principio que torna \emph{individua} a alma humana éa \emph{sua propria essencia}, \emph{emquanto se refere ao corpo}, que ella informa. Na verdade, a nossa alma é uma forma \emph{subsistente} e \emph{incompleta} na \emph{especie}. Por ser forma \emph{subsistente}, a alma torna-se \emph{individua} pela sua propria \emph{essencia}. Por ser \emph{incompleta} na \emph{especie}, a alma, pela sua propria \emph{essencia}, refere-se ao corpo, com que constitue uma substancia completa na sua especie. --- Por outras palavras: A alma é \emph{individua}, emquanto não está dividida em si mesma e está dividida das outras almas. Ora, se a alma não está dividida em si, esta indivisão deriva da sua propria \emph{essencia}, pois ella é essencialmente\pgn{88}\emph{subsistente}; --- se a alma está dividida das outras almas, esta divisão deriva da relação que a alma tem ao corpo, de que é forma. Por isso a \emph{individuação} da alma deriva da sua \emph{essencia}, como do seu \emph{principio} e \emph{causa}, --- e deriva do \emph{corpo}, como do seu \emph{termo}.\par Objectam: A \emph{relação} é uma coisa \emph{accidental}. Ora uma coisa \emph{accidental} não pôde tornar \emph{individua} uma \emph{substancia}, como é a nossa alma; pois a \emph{unidade} identifisa-se com a \emph{entidade}.\par Respondemos que a \emph{relação}, que a alma tem ao seu corpo, não é \emph{categorica}, mas é \emph{transcendental}, e por isso não é uma coisa \emph{accidental}, mas, na \emph{realidade}, identifica-se com a propria \emph{essencia} da alma; pois a alma é ESSENCIALMENTE \emph{acto} do corpo natural e organisado.\par Do que deixamos dicto vê-se claramente que a razão, o fundamento da multiplicidade \emph{numerica} da alma \emph{intellectual} é a sua essencial \emph{relação} á \emph{materia}, que ella informa. Na verdade, a razão ou o fundamento da multiplicidade numerica da alma intellectual deve deduzir-se --- ou da propria alma, --- ou da materia que ella informa. Ora não póde deduzir-se da alma; pois a alma é forma substancial, e a forma substancial é por si \emph{simples} e \emph{indivisivel}, e por isso não tem em si fundamento algum para se multiplicar numericamente na mesma especie. Logo o fundamento da multiplicidade \emph{numerica} da alma intellectual deve deduzir-se necessariamente da materia, que ella informa. Com effeito, a materia, em razão da sua quantidade, pôde dividir-se e effectivamente divide-se em partes, cada uma das quaes, por ter uma determinada extensão, distingue-se de todas as outras, e assim todas formam muitos e diversos sujeitos, \emph{numericamente} distinctos entre si. Ora, se a cada parte da materia se unir uma forma \emph{substancial}, é claro que teremos tantas formas substanciaes, \emph{numericamente} distinctas entre si, quantas forem as partes em que está dividida a materia; porque não pôde multiplicar-se o sujeito, sem que ao mesmo tempo se multiplique a forma. Mas, a alma intellectual está unida ao corpo, como a forma substancial á materia. Logo a alma intellectual multiplica-se \emph{numericamente} na mesma especie pelo facto de se referir a diversos corpos. --- Em poucas palavras: A razão da multiplicidade \emph{numerica} da alma humana não é a \emph{forma}, como \emph{tal}, pois esta é a razão da multiplicidade das \emph{especies} e nao dos \emph{individuos}, --- mas é a \emph{materia determinada pela quantidade}, ou (melhor) é a propria alma emquanto se refere a uma certa e determinada materia. --- Diz S. Thomaz: « Manifestum est, quod de ratione animae humanae est quod corpori humano sit unibilis, cum non habeat in se \emph{speciem completam}, sed speciei complementum sit in hoc composito. Unde quod sit unibilis huic aut illi corpori, multiplicat animam \emph{secundum numerum}, non autem \emph{secundum speciem}; sicut et haec albedo differt ab illa numero per hoc quod est hujus vel illius . subjecti » (\emph{De anima}, a. 9. --- Cf. e. \emph{Gentes}, lib. II, c. 81).\par Pela doutrina exposta podemos julgar da \emph{metempsychose}.\par Este systema, admittido primeiramente no Oriente e depois defendido por \emph{Pythagoras} e \emph{Platão} e por outros, ensinava que as almas transmigram de um para outro corpo. --- Foi resuscitado nos tempos modernos por \emph{Lerou}«, \emph{HReynaud}, \emph{Michelet}.\par A \emph{metempsychose} é falsa. Na verdade, este systema estabelece, como\pgn{89}principio, que a alma é indifferente para informar este ou aquelle corpo.\par \emph{fol} Ora tal principio é falso. Porquanto, cada alma refere-se a um determinado corpo; e esta relação transcendental, que faz com que a alma seja \emph{individua}, singular e diffira \emph{numericamente} das outras almas, constitue a propria essencia da nossa alma. Ora a essencia da alma é immutayvel, e por isso a alma conserva sempre uma relação com o seu corpo, embora esteja separada d'elle. Conservando sempre esta relação, a alma não póde transmigrar de um para outro corpo, pois a forma deve ser proporcionada á materia.\par Advirta-se que a differenca, que existe entre as almas e que é motivada pela sua relação transcendental a diversos corpos, não é uma differenca \emph{especifica}, que constitue as almas em especies distinctas, mas é uma differenga \emph{numerica}, que faz com que muitas almas, pertencentes todas á mesma especie, sejam numericamente distinctas umas das outras.\par Poderiamos tambem mostrar a falsidade da \emph{metempsychose}, atacando o principio, em que se basêa e que é a preexistencia das almas humanas nas estrellas, d' onde teriam sido expulsadas em castigo de culpas commettidas, --- principio que é falso, como demonstraremos mais tarde.\par *\par \emph{A} união da alma com o corpo, \& unidade e a distincção numerica da propria alma são assumptos de tão alta importancia, que a Egreja Catholica, Mestra infallivel de verdade, tem tido o cuidado de as resolver, e a sua resolução está, como não podia deixar de estar, em perfeita conformidade com a razão natural. --- Os principaes documentos da Egreja, a este respeito, são as definições de CLEMENTE V no Concilio de Vienna, de. LrAo X no Concilio \emph{Lateranense} V, e do S. Padre Pro IX nas cartas ao Arcebispo de \emph{Colonia} e ao Bispo de \emph{Breslau}.\par a) \textsc{CLEMENTE} V, no Concilio de Vienna, em 1311, por occasiüo dos erros de Pedro João Oliva, sequaz de Averroes, emanou a seguinte definição: « \emph{Doctrinam omnem seu propositionem}. \emph{temere asserentem}. \emph{vel vertentem in dubium quod}. \emph{substantia}. animae \emph{rationalis seu intellectivae} vere ac \emph{per se humani corporis non sit forma}, \emph{velut erroneam}. ac \emph{catholicae inimicam}. \emph{fidei}, \emph{praedicto sacro approbante Concilio}, \emph{reprobamus}; DEFINIENTES, \emph{ut cunctis nota sit fidei sincera veritas} ac \emph{praecludatur universis erroribus}. \emph{aditus} ne \emph{subintrent}, \emph{quod}. \emph{quisquis deinceps asserere}, \emph{defendere} seu \emph{tenere pertinaciter praesumpserit quod anima rationalis seu intellectiva non sit} FORMA \emph{corporis humuni} PER SE \emph{ef} ESSENTIALITER, \emph{tamguan} HAERETICUS \emph{sit censendus} ».\par b) Lzão X, no Concilio Lateranense V, em 1515, por occasiüo dos erros de Pomponacio, repetiu a mesma definição: « \emph{Hoc sacro approbante Concilio damnamus et reprobamus omnes asserentes animam intellectivam}. \emph{mortalem esse}, \emph{aut unicam}. \emph{in cunctis hominibus}, \emph{et haec in dubium vertentes}: \emph{quum illa non solum} vERE, PER SE \emph{ef} ESSENTIALITER /\emph{vumani corporis} FORMA \emph{existat}, \emph{sicut in canone fel}. \emph{rec}. \emph{Clementis Papae V}, \emph{praedecessoris nostri}, \emph{in generali Viennensi Concilio edito continetur}, \emph{verum et immortalis et pro corporum}, \emph{quibus infunditur}, \emph{multitudine singulariter multiplicabilis et multiplicata}. \emph{et multiplicanda sit} ».\par 0) O S. Padre Prio IX, na sua carta ao Card. \emph{Geissel}, Arcebispo de Co.\pgn{90}lonia, em 1857, condemnando os erros de \emph{Giinther}, que admittia no homem duas almas --- a sensitiva e a intellectual --- e porisso negava que a segunda fosse \emph{forma}, ao menos \emph{immediata}, do corpo, assim se exprime: « \emph{Noscimus iisdem libris} (de Günther) \emph{laedi catholicam sententiam} ac \emph{doctrinam de homine}, \emph{qui corpore et anima ita absolvatur}, \emph{ut anima}, \emph{eaque rationalis}, \emph{Sit} VERA, PER SE, \emph{algue} IMMEDIATA \emph{COTJOris} FORMA 7.\par a) O mesmo S. Padre Pro IX, na sua carta ao Bispo de \emph{Breslau}, om 30 de abril de 1860, depois de declarar que a doutrina de \emph{Daltzer} é similhante á de \emph{Günther}, exalta a doutrina escholastica ácerca da união da alma com O corpo, com as seguintes palavras: » \emph{Hanc sententiam}, \emph{quae unwm in homine ponit vitae principium}, \emph{animam scilicet rationalem}, \emph{a qua corpus quoque et motum et vitam}. \emph{omnem et sensum accipiat}, \emph{in Dei Ecclesia esse communissimam}, \emph{atque doctoribus plerisque et}. \emph{probatissimis quidem maocime}, \emph{cum Ecclesiae dogmate ita videri conjunctam}, \emph{ut hujus sit legitima solaque vera interpretatio}, \emph{nec proinde sine errore in fide possit negari} n.\par Analysemos brevemente estes documentos da Egreja.\par a) A alma humana, segundo o Concilio de Vienna, é \emph{forma} do corpo. Esta palavra \emph{forma} deve tomar-se no sentido em que a empregou o proprio Concilio. Ora o Concilio empregou-a no sentido que lhe dão os escholasticos; pois, em 1811, a palavra \emph{forma} significava uma realidade que aperfeigóa e completa a materia, nem havia divergencia neste ponto.\par \emph{b}) A alma humana não é forma \emph{accidental} do corpo, nom forma \emph{assistente}, mas é forma \emph{substancial}. Na verdade, a alma é forma \emph{per se}, emquanto \emph{naturalmente} tende a communicar o ser ao corpo, --- é forma \emph{essentialiter}, emquanto communiea ao corpo o ser humano pelo facto de se lhe unir na sua \emph{essencia} ou \emph{substancia}, --- é forma \emph{immediata}, emquanto, unindo- 86 ao corpo na sua \emph{essencia}, une-se por si mesma e sem outro meio. Ora a forma, que por si communica á materia o ser, de modo que resulte uma substancia completa, é \emph{substancial}. Logo a alma é forma \emph{substancia}! do corpo.\par c) A alma humana une-se ao corpo quanto á sua \emph{substancia}, ou \emph{essencia}, e não quanto á operação \emph{intellectual}; pois o Concilio de Vienna diz que éa \emph{substancia} da alma racional ou intellectual que se une ao corpo.\par \emph{d}) A alma \emph{intellectual} é á \emph{unica} forma \emph{substancial} existente no homem e é o principio pelo qual o corpo é corpo, é vivo, é sensitivo. D'ahi se segue que, se o corpo fosse corpo independentemente da alma, esta, pelo facto de encontrar o corpo já constituido e por isso dotado de \emph{forma}, não seria forma \emph{substancial} nem \emph{immediata}. --- Por isso todos os que rejeitam o /\emph{ylomorphismo}, e seguem o \emph{atomismo} ou o \emph{dinamismo}, negam ou devem negar, 'se quizerem ser logicos, que a alma é forma \emph{immediata} do nosso corpo.}.

%%>>> art-11

\section*{Artigo III.\\ Séde da alma humana.}
\addcontentsline{toc}{section}{Artigo III. --- Séde da alma humana.}

%%>>> t55

\tese{55}{Estado da questão.}
A alma humana está no corpo, --- não quanto á sua \emph{extensão}, pois ella, por ser simples e espiritual, não é dotada de partes, --- mas quanto á sua \emph{essencia} \pg{87} \pg{88}  e \emph{energia}. --- Agora pergunta-se: Estará a alma em \emph{todo} o corpo e em \emph{cada} parte d'elle? E, no caso d'ella estar em \emph{todo} \pg{89}  o corpo e em \emph{cada} parte do corpo, --- estará \emph{toda} em \emph{todo} o corpo e \emph{toda} em \emph{cada} parte d'elle, --- e só quanto á totalidade \pg{90}  da \emph{essencia}, ou tambem quanto á totalidade da \emph{energia}? --- A estas perguntas respondemos com as seguintes proposições \footnote{A questão Acerca da \emph{séde da alma humana} é um simples corollario da outra questão ácerca da \emph{união da alma com o corpo}. Se a alma humana fosse apenas uma forma \emph{accidental} ou \emph{assistente} do nosso corpo, poderia residir numa determinada parte do organismo. Mas, sendo ella uma forma \emph{substancial}, deve residir \emph{toda} em \emph{todo} o corpo, \emph{e toda} em \emph{cada} parte d'elle.\par São duas, pois, as opiniões ácerca da \emph{séde} da alma. A primeira admitte que a alma reside numa certa o determinada parte do organismo, --- a segunda sustenta que existe \emph{toda} em \emph{todo} o corpo e \emph{toda} em \emph{cada} parte d'elle.\pgn{91}--- No texto defendemos a segunda opinião, que é a de \emph{Aristoteles}, de Santo \emph{Agostinho}, de Santo \emph{Thomaz} e de todos os \emph{escholasticos}; e no fim do presente artigo, em nota, combateremos a primeira.\par S. Thomaz trata d'esta questão na \emph{Sum}. \emph{Theol}., p. l, q. 76, a. 8; --- \emph{contr}. \emph{Gent}., l. IL, c. (9; --- \emph{De anima}, a. 10; --- \emph{De spirit}. \emph{creat}., a. 4.\par Para a exacta intelligencia d'esta questão, devemos notar duas coisas:\par a) Uma substancia \emph{creada} pode estar num logar por dois modos: \emph{circumscriptivo} e \emph{definitivo}. --- Está por um modo \emph{cirewnscriptivo}, quando uma parte da substancia está numa parte do logar e toda a substancia em todo o logar, e assim a substancia é contida e medida pelo logar. Este modo é proprio das substancias \emph{corporeas}. --- Está por um modo \emph{definitivo}, quando a substancia, embora só se encontre num determinado logar, todavia está toda em todo o logar e toda em cada parte do logar, e assim não é contida pelo logar, mas antes é ella que contém o logar. Este modo é proprio das substancias \emph{espirituaes} (Cf. \emph{Ontologia}, cap. IV, art. \emph{V}). --- Dizemos --- uma substancia \emph{creada}, porque a substancia \emph{increada} occupa todos os logares, a todos se extende, e a todos excede pela sua immensidade.\par b) Ha uma triplice \emph{totalidade}, conforme a triplice diversidade de partes: isto é, \emph{quantitativa}, \emph{essencial} e \emph{virtual}. A totalidade \emph{quantitativa} é aquella pela qual o todo se diz \emph{quanto} e por isso se compõe de partes \emph{quantitativas}. Esta é propria da materia, e sô \emph{accidentalmente} pode convir a algumas formas, que se extendem por cima da quantidade e que por isso se dividem pela divisão da quantidade: assim divide-se a \emph{alvura}, quando se divide a superficie alva. --- A totalidade \emph{essencial} diz respeito á \emph{essencia}, e por isso se compõe de partes \emph{essenciaes}, que podem ser \emph{intencionaes} (como o \emph{genero} e a \emph{differença}) ou \emph{reaes} (como á \emph{materia} e a \emph{forma}). --- A totalidade. \emph{virtual} refere-re \& \emph{virtude} ou \emph{energia}, e por isso se compõe de partes \emph{viriuaes} ou \emph{potestativas}. --- Não convindo a totalidade \emph{quantitativa} á alma \emph{humana}, que nem \emph{por si} nem \emph{accidentalmente} se pode dividir, a questão limita-se ás outras duas totalidades, isto é, se a alma está no corpo quanto á totalidade da \emph{essencia} e da \emph{energia}.}.

%%>>> t56

\pg{91} \tese{56}{A alma esta em todo o corpo e em cada parte d'elle.}

a) Está em \emph{todo} o corpo. --- Na verdade, a alma, como provámos, é a forma \emph{substancial}, que constitue o nosso corpo numa certa e determinada especie, communieando-Ihe a vida vegetativa e sensitiva. Ora a forma \emph{substancial} deve estar na potencia, que ella determina e á qual communiea certas propriedades, e \emph{todo} o nosso corpo é humano e dotado de vida vegetativa e sensitiva. Logo a alma está em \emph{todo} o corpo.

b) Está em \emph{cada} parte do corpo. --- Com effeito, cada parte do nosso corpo pertence á especie humana e possue vida vegetativa e sensitiva; de modo que a alma é forma \emph{substancial} não só de \emph{todo} o corpo, mas tambem de \emph{cada} sua parte. Ora a forma está no ente de que é forma. Logo a alma está em \emph{cada} parte do nosso corpo, --- De mais, se a alma estivesse só em alguns orgãos do corpo, os outros orgãos deveriam ser \pg{92}  aperfeiçoados por outras formas \emph{substanciaes}, e assim haveria no homem varias formas \emph{substanciaes}, o que é impossivel \footnote{S. Thomaz diz: « Totum corpus et omnes ejus partes habent esse \emph{substantiale} et \emph{specificwm} per animam...Sic igitur, cum omnis actus sit in eo cujus est actus, oportet animam, quae est actus \emph{totius} corporis et \emph{omnium partium}, essein toto corpore etin qualibet ejus parte » (\emph{De sp}. \emph{creat}., a. 4).\par Objectam: Se a nossa alma estivesse em \emph{cada} parte do corpo, estaria em varios logares. Ora um ente não póde estar em varios logares.\par Resposta. E' certo que um ente \emph{material} não póde estar, ao mesmo tempo, em varios logares, pois que é limitado e circumseripto pelo logar em que está; mas não repugna que um ente \emph{espiritual}, que não é circeumscripto nem limitado pelo logar, se encontre em cada parte do corpo extenso, vivificando-o e constituindo-o numa certa e determinada especie. --- Além d'isso, é falso que a alma estaria ao mesmo tempo em varios logares, se residisse em cada parte do corpo. Com effeito, a alma está no corpo, não como uma substancia \emph{completa} está num logar, mas como a forma está na materia; e as partes que ella informa não estão separadas, mas estão unidas de maneira que constituem um todo completo.}.

%%>>> t57

\tese{57}{A alma, quanto á essencia, está toda em -todo o corpo e toda em cada parte d'elle.}
A \emph{essencia} de uma substancia espiritual é \emph{simplicissima}, e por isso não póde divider-se em diversas partes, que oceupem diversos logares; mas, onde está, está \emph{toda}. A \emph{essencia} da alma humana, que é substancia espiritual, é \emph{simplicissima}. Logo, onde está, está \emph{toda}. Ora a essencia da alma está em \emph{todo} o corpo e em \emph{cada} parte d'elle. Logo \emph{toda} a essencia da alma, isto é, a alma, quanto á \emph{essencia}, está \emph{toda} em \emph{todo} o corpo e \emph{toda} em \emph{cada} parte d'elle \footnote{Objectam: Se a alma estivesse \emph{toda} em todas as partes do organismo, augmentando oudiminuindo estas, deveria augmentar ou diminuir aquella. Ora a alma, por ser simples, é incapaz de augmento e de diminuição.\par Respondemos que, augmentando ou diminuindo as partes do organismo, a alma exerce a sua actividade numa quantidade de materia mais ou menos extensa, começa a vivificar novas partes. que antes não vivifieava, ou deixa de vivifiear as que vivificava, sem que por isso ella augmente ou diminua; assim como uma luz não augmenta nem diminue, pelo facto de illuminar um quarto mais ou menos extenso. Em resumo: o augmento e a diminuição dão-se no \emph{termo}, que é o corpo, e não no \emph{principio}, que é a alma.}.

%%>>> t58

\tese{58}{A alma, quanto á energia, não esta toda em todo o corpo nem toda em cada parte d'elle.}
Na verdade, a energia da alma é multipla e diversa, porque constitue a vida vegetativa, a sensitiva e a intellectual. Ora, se a energia que constitue a vida \emph{intellectual} se exerce sem orgão, a que constitue a vida \emph{vegetativa} e a \emph{sensitiva} precisa do auxilio dos orgãos. E, como os orgãos não existem em \emph{todo} o corpo nem em \emph{cada} parte d'elle, a energia que constitue a vida vegetativa e a \pg{93}  sensitiva n\&o existe em todo corpo, nem em cada parte d'elle, mas só em \emph{algumas} regiões. Logo a alma, quanto á \emph{energia}, não, está \emph{toda} em todo o corpo, nem \emph{toda} em cada parte d'elle.

As regiões do corpo humano, ás quaes a alma communiea maior energia, são o \emph{cerebro} e o \emph{coração} ; pois este é o orgão principal da vida \emph{vegetativa}, a aquelle o é da vida \emph{sensitiva} \footnote{Não repugna que á \emph{essencia} da alma resida \emph{toda} em \emph{toda} o corpo e \emph{toda} em \emph{cada} parte d'elle, \emph{e} que no emtanto a sua \emph{energia} esteja limidada a determinados pontos do organismo. Por quanto, sendo a energia ou actividade da alma realmente distincta da sua essencia, como provaremos d'aqui a pouco, segue-se que á alma póde estar em todo o corpo com \emph{toda} \& sua essencia, sem que ahi esteja com \emph{toda} a sua energia.\par Em quanto á vida \emph{intellectual}, isto é, emquanto ás faculdades \emph{intellectuaes}, a alma não está em \emph{nenhuma} parte do organismo, porque essas faculdades, como vimos, não são \emph{organicas}, mas \emph{espirituaes}. --- Todavia nada obsta a que, em harmonia com o senso commum, se colloque a intelligencia na cabeça ou no cerebro. A razão d'isto é que as operações intellectuaes, embora não dependam do orgão material, comtudo não se podem realizar sem o concurso da imaginação, e esta reside numa determinada parte do cerebro: d'ahi a fadiga e o cançaço que sentimos no cerebro, durante uma intensa applicação da intelligencia.\par S. Thomaz resume toda esta questão nas seguintes palavras: « Anima non habet extensionem in materia: unde in ea prima totalitas (secundum \emph{quantitetem}) locum non habet. Relinquitur ergo quod \emph{secundwm totalitatem essentiae} simpliciter enuntiari possit esse tota in qualibet corporis parte, non \emph{autem}. \emph{secundum totalitatem virtutis}; quia partes difformiter perficiuntur ab ipsa ad diversas operationes; et aliqua operatio est ejus, scilicet \emph{intelligere}, quam per \emph{nullam} partem corporis exequitur. Unde sic accepta totalitate animae \emph{secundum}. \emph{virtutem}, non solum non est tota in qualibet parte, sed nec est tota in toto: quia virtus animae capacitatem corporis excedit » (\emph{De spirit}. \emph{creat}., a. 4).\par Examinemos a opinião dos que, admittindo uma união \emph{accidental} entre a alma e o corpo, limitam a alma a um ponto determinado do corpo.\par \emph{Platão} collocou a alma humana na \emph{cabeça}; \emph{Parmenides}, no \emph{peito}; \emph{Epicuro} e os \emph{Estoicos}, no \emph{coração}; \emph{Empedocles}, na \emph{substancia do sangue}; \emph{Pyhagoras}, no \emph{coração} e na \emph{cabeça}; \emph{Descartes} e \emph{Mallebranche}, no \emph{cerebro}; \emph{La Peyronnie}, no \emph{corpo calloso}; e outros, nas \emph{meninges}, no \emph{cerebello}, ete.\par Estas opiniões são falsas pelos seguintes fundamentos:\par a) Baseam-se num falso supposto; porque partem do principio de que a união entre a alma e o corpo é apenas \emph{accidental}. Ora este principio, como demonstr\&àmos, é falso.\par \emph{b}) São contrarias á experiencia. Na verdade, a experiencia attesta que \emph{todas} e \emph{cada} uma das partes do nosso organismo são dotadas de vida vegetativa e sensitiva. Ora seria impossivel que todas e cada uma das partes do nosso organismo fossem dotadas de vida vegetativa e sensitiva, se a alma, que é o principio d'aquella duplice vida, residisse apenas numa certa e determinada região do corpo humano.\par c) Confundem o \emph{indivisivel material} com o \emph{indivisivel espiritual}, E' certo\pgn{94}que o \emph{indivisivel material}, por ser o termo da extensão, não póde estar, ao mesmo tempo, em varias partes do corpo extenso, mas deve occupar um ponto certo e determinado. Mas o \emph{indivisivel espiritual} não é extenso, nem é o termo da extensão, e por isso não se segue que só possa e deva residir num certo e determinado ponto do corpo humano.}.

\textbf{FACULDADES DO HOMEM}

%%>>> t59

\pg{94} \tese{59}{Divisão da secção segunda.}
Da essencia emanam as faculdades. Por isso, tendo considerado a \emph{essencia} do homem, devemos agora considerar as suas \emph{faculdades}. --- Dividimos a secção \emph{segunda} em \emph{tres} capitulos. No \emph{primeiro} capitulo tratamos das faculdades humanas \emph{em geral}; no \emph{segundo}, das faculdades \emph{sensitivas} do homem; no \emph{terceiro}, das suas faculdades \emph{intellectuaes}.

%%>>> cap-2

\chapter*{CAPITULO PRIMEIRO\\[0.3em]\large FACULDADES HUMANAS EM GERAL}
\addcontentsline{toc}{chapter}{CAPITULO PRIMEIRO --- FACULDADES HUMANAS EM GERAL}

{\small\noindent\textsc{Summario}: --- As faculdades consideradas em si mesmas. --- As faculdades consideradas em relação á alma. --- As faculdades consideradas em re- lação ás suas operações.\par}

%%>>> art-12

\section*{Artigo I.\\ As faculdades consideradas em si mesmas.}
\addcontentsline{toc}{section}{Artigo I. --- As faculdades consideradas em si mesmas.}

%%>>> t60

\tese{60}{Faculdade.}
\emph{Faculdade}, em geral, é o \emph{principio proximo} e \emph{immediato} das nossas operações. --- O principio de que emanam as operações é a propria \emph{forma} do ente. Mas a \emph{forma} do ente não opéra \emph{immediatamente por si}, mas por meio de certas \emph{qualidades}, que lhe são inherentes. Estas \emph{qualidades} denominam-se \emph{faculdades} ou \emph{potencias}. --- Devemos, pois, distinguir dois principios de operações: um \emph{mediato} e \emph{remoto}, que é a \emph{alma}; outro \emph{immediato} e \emph{proximo}, que é a \emph{faculdade}.

Neste artigo consideramos as faculdades humanas \emph{em si mesmas}; e nos dois artigos seguintes consideramol-as em relação ao \emph{principio}, de que derivam e que é a \emph{alma}, e em relação ao \emph{termo}, para o qual tendem e que consiste nas \emph{operações} \footnote{O tratado das \emph{faculdades} costuma chamar-se \emph{dynamilogía}.}.

%%>>> t61

\pg{95} \tese{61}{As faculdades humanas são muitas e diversas.}
O homem é o \emph{animal racional}. Logo deve ser dotado de todas as faculdades, que possuem os animaes e os entes dotados de razão. Ora as faculdades, que possuem os animaes e os entes dotados de razão, são muitas e diversas. Logo as faculdades humanas são muitas e diversas \footnote{Diz S. Thomaz: » Est et alia ratio quare anima humana abundat diversitate potentiarum, videlicet quia est in confinio spiritualium et corporalium ereaturarum; et ideo concurrunt in ipsa virtutes utrarumque creaturarum » (\emph{Sum}. \emph{th}, p. I, q. 77, a. 2).\par A pluralidade das faculdades não é contraria á \emph{simplicidade} da alma. Com effeito, a alma, relativamente ás suas faculdades, não é um \emph{todo éntegrante}, isto é, não é um todo composto de partes separaveis, como é uma casa; mas é um todo \emph{potencial}, isto é, é um todo que, embora seja simplicissimo, todavia é dotado de varias e diversas faculdades, de que se serve, como de instrumentos, para exercer a sua multipla e variada energia.}.

%%>>> t69

\tese{69}{Divisão das faculdades humanas.}
As faculdades humanos dividem-se, principalmente, --- a) em \emph{vegetativas}, \emph{sensitivas} e \emph{intellectuaes}, --- \emph{b}) em \emph{activas} e \emph{passivas},

a) \emph{Faculdades vegetativas}, \emph{sensitivas} e \emph{intellectuaes}. - As classes supremas das faculdades humanas são tantas, quantas são as classes supremas das operações humanas; porque as faculdades são o principio proximo e immediato das nossas operações. Ora são tres as classes supremas das operações humanas, porque estas pertencem ou á vida \emph{vegetativa}, ou á vida \emph{sensitiva}, ou \& vida \emph{intellectual}. Logo são tambem tres as classes supremas das nossas faculdades: \emph{vegetativas}, \emph{sensitivas} e \emph{intellectuaes}. --- As faculdades \emph{vegetativas} e \emph{sensitivas} chamam-se \emph{organicas}, porque residem nos orgãos vivifieados pela alma; as \emph{intellectuaes} dizem-se \emph{inorganicas} ou \emph{espirituaes}, porque não residem nos orgãos animados, mas exclusivamente na alma \footnote{Esta classificação das faculdades deve preferir-se ás outras, porque é \emph{adequada}, \emph{irreductivel} e \emph{ordenada}. --- E' \emph{adequada}, porque não ha operação \emph{especifica}, que não seja produzida por uma das faculdades enumeradas. --- T? \emph{irreductivel}, porque as operações relativas ás faculdades de uma especie não podem ser produzidas pelas faculdades de outra especie. --- E' \emph{ordenada}, porquanto, sendo uma faculdade mais ou menos nobre conforme o seu objecto fór mais ou menos elevado, segue-se que --- as faculdades \emph{vegetativas}, que tem por objecto o corpo organico unido á alma, occupem o infimo grau, --- as faculdades \emph{sensitivas}, que tem por objecto os entes sensiveis e tendem para elles, occupem um grau superior ás faculdades vegetativas, --- e as faculdades \emph{intellectuaes}, que têm por objecto o ente \emph{immaterial}, occupem o mais alto grau (\emph{Of}. \emph{Sum}. \emph{Theol}. p. 1, q. 78, a. 1).\par As opiniões dos philosophos modernos ácerca da classificação das faculdades\pgn{96} humanas são muitas e diversas, mas podem reduzir-se a duas principaes, que hoje são mais seguidas.\par Uns dividem as faculdades humanas em \emph{sensibilidade}, \emph{intelligencia} e \emph{vontade}. --- Esta classificação é imperfeita, não só porque não enumera as faculdades \emph{vegetativas}, mas tambem porque exprime as \emph{especies} das faculdades \emph{intellectuaes} e apenas denota o \emph{genero} das faculdades \emph{sensitivas}.\par Outros dividem as faculdades humanas em \emph{sensibilidade}, \emph{intelligencia} e \emph{actividade}. --- Esta classificação tambem é imperfeita, não só porque não enumera as faculdades \emph{vegetativas}, mas tambem porque admitte a \emph{actividade} como uma faculdade especial, quando é certo que tanto a \emph{sensibilidade} como a \emph{intelligencia} são dotadas de actividade.}.

b) \emph{Faculdades activas} e \emph{passivas}. - Todas as faculdades, que enumerámos, ou são \emph{activas}, ou são \emph{passivas}. - As faculdades são \emph{activas}, quando transmutam ou produzem o objecto, fazendo com que este de \emph{potencial} se torne \emph{actual}; taes são as faculdades \emph{vegetativas}, que transformam o alimento e o convertem na substancia do ser vivo, e tal é a intelligencia \emph{agente}, que, como veremos, torna intelligivel \emph{em acto} o objecto que só o era \emph{em potencia}. - As faculdades são \emph{passivas}, quando são transmutadas ou determinadas ao acto pelo objecto já existente; taes são as faculdades \emph{sensitivas}, que são movidas ao acto da sensação pelos proprios objectos sensiveis, e taes são a \emph{tintelligencia possivel}, que é determinada ao acto de entender pela imagem espiritual do objecto, e a \emph{vontade}, que tambem é movida pelo bem preexistente e conhecido. --- Por isso o \emph{objecto} --- é \emph{paciente} e \emph{termo} com relação ás faculdades \emph{activas}, porque é elle que, sob acção da faculdade, se muda e passa da potencia ao acto, --- mas é \emph{agente} e \emph{principio} com relação ás faculdades \emph{passivas}, porque é sob a acção d'elle que a faculdade se muda e passa da potencia ao acto \footnote{A definição das faculdades \emph{activas} e \emph{passivas}, apresentada no texto, indica claramente que consideramos as mesmas faculdades em ordem ao \emph{objecto}, e não em ordem á \emph{operação}. Por quanto, se a faculdade se considera em ordem á \emph{operação}, então deve dizer-se que toda a faculdade \emph{vital} é, ao mesmo tempo, \emph{activa} e \emph{passiva}; porque toda a faculdade \emph{opera} (e por isso pode chamar-se \emph{activa}), \emph{e} na sua operação é sempre \emph{movida} por algum agente externo (e neste sentido á \emph{passiva}). Mas, se á faculdade se considera em ordem ao \emph{objecto}, então é \emph{activa} se transmuta o objecto, e é \emph{passiva} se é transmutada pelo objecto. --- Erra, pois, \emph{Galluppi} (\emph{Lesioni di Filosofia}, lez. 96), quando ensina que a faculdade \emph{activa} é a que pode fazer alguma coisa, e faculdade \emph{passiva} é á que só é .capaz de receber em si a acção d'um outro agente. Faculdade que nada pode fazer não é, nem se pode chamar faculdade, porque faculdade é essencialmente um principio de operações. (Cf. S. Th., \emph{De veritate}, q. 16, a. 1 ad 13). ' Com as faculdades \emph{passivas} não se deve confundir a faculdade, que os\pgn{97}theologos chamam \emph{obediencia!l}, \emph{e} que consiste na capacidade ou aptidão que a creatura tem para receber de Deus a virtude ou a energia de produzir effeitos superiores ás forças naturaes. Nem isto repugna. Como as creaturas dependem, no seu ser e nas suas operações, da Causa primeira, que é Deus, pode Elle elevar a energia das creaturas a uma ordem sobrenatural, de modo que possam. produzir actos superiores ás forças naturaes. (Cf. S. Thom. \emph{Sum}. \emph{Theol}, p. 9, q. 11, a. 1; - \emph{De potentia}, q. 6, a. 1 ad 18).}.

63. As faculdades conhecem-se \textbf{pelas operações} e \textbf{pelos objectos}.

a) As \emph{faculdades conhecem-se pelas operações}. --- Na verdade, todas as causas conhecem-se pelos effeitos, que ellas devem naturalmente produzir. Ora as operações são effeitos, que as faculdades devem naturalmente produzir. Logo as faculdades conhecem-se pelas \emph{operações}.

b) As \emph{faculdades conhecem-se pelos objectos}. --- As faculdades, como provámos, conhecem-se pelas suas operações. Mas as operações conhecem-se pelos seus \emph{objectos}. Com effeito, uma coisa conhece-se por outra, quando esta é conhecida primeiro do que aquella; pois do conhecido subimos ao desconhecido. Ora os objectos, que são o termo attingido pelas operações, são conhecidos anteriormente ás proprias operações; pois que a intelligencia \emph{primeira} e \emph{directamente} apprehende o objecto, e em seguida, reflectindo sobre si mesma, apprehende a operação. Logo as operações e porisso as faculdades conhecem-se pelos \emph{objectos} \footnote{\emph{Rosmini} sustenta que o meio, que nos leva ao conhecimento das faculdades, é a essencia da alma, e não as operações nem os objectos.\par E' necessario distinguir. Ou trata-se de um conhecimento \emph{scientifico} das faculdades, ou apenas de um conhecimento \emph{inicial}. --- Sese trata de um conhecimento \emph{scientifico} das faculdades, concordamos em que o meio, que nos leva a esse conhecimento, é a propria \emph{essencia} da alma; porque não se conhece \emph{scientificamente} uma coisa, senão por meio da sua \emph{causa}, e a causa, de que emanam as faculdades, é a essencia. --- Se, porém, se trata de um conhecimento 2;5\emph{icial}, então o meio, que nos leva a esse conhecimento, consiste nos actos que as faculdades produzem, e nos objectos que ellas attingem.\par A verdadeira doutrina, pois, àcereg do conhecimento das faculdades á a seguinte: Primeiramente, de um imperfeito conhecimento das operações e das faculdades subimos ao conhecimento da essencia da alma; em seguida, do conhecimento da essencia da alma descemos ao conhecimento das faculdades e das operações. Cf. S. Th., \emph{De veritate}, q. X, a. 9.\par 7 Vor. II.}.

64. As faculdades especificam-se pelas operações e pelos objectos.

a) As \emph{faculdades especificam-se pelas suas operações}. --- Os entes que tendem naturalmente para certos e determinados fins ou termos, especifieam-se pelos mesmos fins ou termos; porque essa tendencia natural exige e suppõe uma proporção \pg{98}  entre o ente e o proprio fim ou termo. Ora as faculdades têm uma tendencia natural para as suas operações. Logo as faculdades especificam-se pelas suas operações. - Por isso uma operação \emph{espiritual} exige e suppõe uma faculdade \emph{espiritual}.

Todavia, as faculdades não desumem a propria especie de toda e qualquer operação, mas só das operações \emph{especificamente} diversas entre si; assim \emph{entender} e \emph{querer}, por serem operações \emph{especificamente} diversas, exigem duas faculdades \emph{especificamente} distinctas, ao passo que \emph{camiünAar} e \emph{correr} não suppõem faculdades \emph{especificamente} distinctas, visto que o \emph{mais} e o \emph{menos} não mudam a especie dos entes.

b) As \emph{faculdades especificam-se pelos seus objectos}. --- As faculdades, como provámos, especifieam-se pelas suas operações. Mas as operações especifieam-se pelos seus \emph{objectos}. Por quanto, as operações derivam ou das faculdades \emph{activas}, ou das faculdades \emph{passivas}. Se derivam das faculdades \emph{activas}, as operações referem-se aos seus \emph{objectos}, como a seus \emph{fins} ou \emph{termos}, e d" elles desumem a propria especie; porque o \emph{fim} é a razão pela qual o ente possue uma determinada natureza, devendo a natureza ser adaptada ao fim. Se derivam das faculdades \emph{passivas}, as operações referem-se aos seus objectos, como ás suas \emph{causas} ou \emph{principios activos}, e d'elles tambem desumem a especie; porque o \emph{agente} imprime sempre a propria \emph{similhanca} no effeito. Logo as operações e porisso as faculdades especifieam-se pelos \emph{objectos}. --- Logo:objectos \emph{espirituaes} exigem e suppoem operações e faculdades \emph{espirituaes}.

Todavia as operações, como tambem as faculdades, não desumem a propria especie dos objectos \emph{materiaes}, isto é, das coisas, que as faculdades attingem, --- mas sim dos objectos \emph{formaes}, isto é, das \emph{propriedades}, que existem nas coisas,e constituem o objecto \emph{proprio} e \emph{directo} das mesmas faculdades \footnote{Esta conclusão é um corollario da conclusão antecedente. Se as operações e os objectos são o principio ou meio pelo qual conhecemos as varias faculdades, segue-se que a especificação tem o sen fundamento nas proprias operações que ellas produzem e nos objectos que attingem.\par Desenvolvamos este ponto que é da maxima importancia pelas suas applicagóes. --- O objecto de uma faculdade á \emph{material} ou \emph{formal}. Objecto \emph{material} da faculdade é a \emph{coisa}, que \& propria faculdade attinge; assim uma \emph{pedra} é o objecto \emph{material} da \emph{vista}. Objecto \emph{formal} é a \emph{propriedade}, que o objecto \emph{material} possue e por meio da qual é attingido pela faculdade; assim a cór da pedra é o objecto \emph{formal} da \emph{vista}, porque a faculdade \emph{visual} não poderia perceber a pedra, se esta não tivesse cór,\par Ora a especificação, e por isso a unidade ou multiplicidade das facul\pgn{99}dades não se deduz da unidade ou multiplicidade do objecto \emph{material}, mas da unidade ou multiplicidade do objecto \emph{formal}. Por quanto, muitos objectos \emph{materiaes} podem referir-se a \emph{uma} só faculdade, e constituir \emph{um} só objecto \emph{formal}, emquanto todos são attingidos pela faculdade sob o mesmo \emph{aspecto}, ou por meio de uma e a mesma propriedade; --- como, pelo contrario, um só objecto \emph{material} pode referir-se a \emph{varias} faculdades, e assim tornar-se \emph{formalmente} multiplo, emquanto pôde ser considerado sob diversos \emph{aspectos}, ou nas suas differentes \emph{propriedades}, cada uma das quaes se refere a uma faculdade especial. Assim o sol, a lua, as plantas, os homens, são objectos \emph{materialmente} diversos, e todavia, para serem percebidos, não exigem diversas faculdades, mas basta uma só: \& \emph{vista}; pois que todos aquelles objectos, \emph{materialmente} diversos, constituem um só objecto \emph{formal}, sendo todos attingidos pela faculdade \emph{visual}, emquanto são dotados de \emph{cór}, que é o objecto \emph{proprio da vista}. Pelo contrario, uma \emph{laranja} póde referir-se a varias faculdades --- 4 \emph{vista}, ao \emph{gosto}, ao \emph{tacto}; porque póde ser attingida por diversas faculdades segundo a diversidade das suas propriedades, que são a cór, o \emph{sabor}, \& \emph{dureza}, etc., e que são irreductiveis á mesma faculdade.\par S. Thomaz trata d'esta questão sobretudo-na \emph{Sum}. \emph{Theol}., p. l, q. \emph{V}7, a. 9, --- nas \emph{Quest}. \emph{Disp}. - \emph{De Anima}, a. 189, - \emph{De veritate}, q. 15, a. 2.\par Para o exacto entendimento d'esta verdade, notamos que ha coisas \emph{absolutas} e coisas \emph{relativas}. Assim \emph{a substancia} é uma coisa \emph{absoluta}. porque se nao refere a outra, mas tem o seu termo em si mesma. Mas não se pode dizer o mesmo das \emph{faculdades}. Estas são os meios de que a substancia se serve para produzir certas operações e attingir certos objectos; por isso são essencialmente \emph{relativas}, isto é, a sua essencia consiste na propria relação ao acto e, por este, ao objecto: \emph{potentia est nihil aliud quam quidam ordo ad actum} (S. Th. \emph{De Anima}, lib. 2, lect. 11). Portanto, se o ser de uma coisa \emph{absoluta} é especificado pela propria \emph{forma}, o ser de uma coisa \emph{relativa}, como é a \emph{faculdade}, deve ser especificado pela outra coisa, á qual ella essencialmente se refere; de modo que é da natureza da \emph{operação} e do \emph{objecto} que depende e se deduz a natureza ou \& especie da \emph{fuculdade}. E como a unidade e a distinegão seguem o ser \emph{especifico} do ente, é claro que d' aquella coisa, de que deriva para o ente o ser \emph{especifico}, deriva tambem a unidade e a distincção. Logo concluimos que das operações e dos objectos deriva para as faculdades não só o ser \emph{especifico}, mas tambem a unidade e a distineção; de modo que serão tantas as faculdades humanas, quantas forem as operações \emph{especificamente} diversas que o homem exerce e quantos os objectos \emph{formaes} das mesmas operações. --- Advirta-se que a \emph{operação é o fim tintrinseco} da faculdade, e o \emph{objecto} é o fim \emph{extrinseco}.\par Objectam: A \emph{causa} não pode ser especificada pelo \emph{e/feito}, porque a causa \emph{precede} e \emph{produz} o effeito, \emph{e} o ente que especifica deve ser anterior ao ente que é especificado. Ora as operações são \emph{e/feitos} das faculdades, e por isso não são as operações que especificam as faculdades, mas são as faculdades que especificam as operações.\par Resposta. Concedemos que a causa não pode ser especificada pelo effeito emquanto \emph{tal}, isto a, emquanto é \emph{e/feito}; mas negamos que a causa não pode\pgn{100}ser especificada pelo effeito emquanto este é o \emph{fim} ou \emph{termo}, para o qual a causa tem uma tendencia ou ordem natural. Ora a operação é o \emph{fim} ou \emph{termo} da faculdade. Por quanto á operação de uma faculdade pode considerar-se --- na ordem da \emph{execução} e na ordem da \emph{intenção}. Considerada na ordem da \emph{execução}, a operação é um \emph{effeito} da \emph{faculdade}, e, como tal, não \emph{especifica} a faculdade, mas apenas a \emph{manifesta}, ou a torna \emph{conhecida}; assim operações espirituaes diversas suppõem e manifestam faculdades espirituaes diversas. Considerada na ordem da \emph{intenção}, isto é, na mente do Creador, que em vista de uma operação cria uma faculdade proporcionada, neste caso a operação é o \emph{fim} ou \emph{termo} da faculdade, e, emquanto tal, é a razão do ser e da natureza d'ella, e por isso a \emph{especifica} e \emph{precede}.\par Insistem: O que especifica uma coisa deve ser \emph{intrinseco} á propria coisa. Ora os objectos são \emph{ewtrinsecos} á faculdade. Logo não a podem especificar.\par Resposta. O que especifica uma coisa deve ser \emph{intrinseco} á propria coisa, quando elle é \emph{parte integrante}, mas não quando é causa ou \emph{termo}. Ora o objecto especifica a faculdade não porque seja \emph{parte integrante} da mesma faculdade, mas porque é \emph{causa} (da faculdade \emph{passiva}) ou é \emph{termo} (da faculdade \emph{activa}). Logo o objecto, embora seja \emph{extrinseco} á faculdade, todavia especifica a propria faculdade pelo modo exposto. --- O principio \emph{intrinseco}, que especifica a faculdade, é a propria disposição ou adaptação intrinseca da faculdade para a sua operação e para o seu objecto. Esta disposição ou adaptação é á propria entidade da faculdade, feita e determinada segundo a natureza e as exigencias da operação.\par S. Thomaz responde a estas duas difficuldades com as seguintes palavras. « Actus, licet sit \emph{posterior} potentia in \emph{esse}, est tamen \emph{prior} in \emph{intentione} et secundum \emph{rationem}, sicut finis in agente. Objectum autem, licet sit \emph{extrinsecuim}, est tamen \emph{principium} vel \emph{finis} actionis. \emph{Principio} autem et \emph{fini proportionantur} ea quae sunt \emph{intrinseca} rei » (\emph{Sum}. \emph{Theol}., p. 1, q. \emph{Vt}, a. 8 ad 1). E' uma synthese clara e profunda de toda a questão em poucas palavras. Note-se sobretudo a ultima maxima, isto é, que a constituição intrinseca de uma coisa é proporcionada ao seu principio e ao seu fim.}.

65. \textbf{Ordem} nas \textbf{faculdades} humanas. --- Nas muitas e diversas faculdades humanas existe uma triplice ordem: de \emph{geração}, de \emph{dependencia}, de \emph{dignidade}. --- Se se considera a \emph{geração}, as \pg{100}  faculdades da vida inferior precedem as da vida superior, porque a natureza vai do imperfeito para o perfeito (assim as faculdades \emph{vegetativas} precedem as \emph{sensitivas}, e estas as \emph{intellectuaes}) --- Se se considera a \emph{dependencia}, as faculdades inferiores dependem das superiores, emquanto estas dirigem e regulam aquellas (assim o vontade diz-se a rainha das faculdades), e, pela sua vez, as superiores dependem das inferiores, emquanto o exercicio d'aquellas suppõe e exige o exercicio d' estas (assim a intelligencia depende dos sentidos). --- Se se considera a \emph{dignidade}, as faculdades que têm por objecto o ente universal são mais nobres do que as que se referem a um ente particular (assim a intelligencia é mais nobre do que a sensibilidade) \footnote{Não se passa da unidade para a multidão sem uma lei de ordem. Por isso deve haver e ha ordem nas muitas e diversas faculdades, que derivam\pgn{101} da unica essencia da alma. --- Esta ordem funda-se --- na propria alma, que, embora una na essencia, tem comtudo aptidão para muitos e dirsos actos, --- nas \emph{faculdades}, das quaes uma depende da outra, --- nas \emph{eracües}, das quaes uma serve a outra, --- na qualidade dos \emph{objectos}. --- uando nas faculdades a ordem se considera só em relação aos \emph{objectos}, do etc de um objecto ser mais nobre do que um outro não se segue necessaamente que uma faculdade dependa da outra no desenvolvimento da sua zvida (Cf. S. Th. \emph{Sum}. \emph{T}., p. I, q. 7, a. 4).}.

%%>>> art-13

\section*{Artigo II.\\ As faculdades consideradas em relação á alma.}
\addcontentsline{toc}{section}{Artigo II. --- As faculdades consideradas em relação á alma.}

\pg{96} \pg{97} \pg{99} \pg{101} A

66. A alma é o principio de todas as faculdades humanas. --- \emph{Principio} de uma faculdade é o ente, de que, como de uma raiz, deriva a propria faculdade. Ora a alma é o ente, de que derivam \emph{todas} as faculdades humanas. Na verdade, todas as nossas faculdades são \emph{tnorganicas} ou \emph{organicas}. Se são \emph{norganicas}, derivam da alma; porque o inorganico ou espiritual só póde resultar de um ente inorganico ou espiritual, como é á nossa alma. Se são \emph{organicas}, tambem derivam da alma; porque, embora residam no corpo, como em seu \emph{sujeito}, todavia residem nelle, não emquanto é corpo (aliás todo o corpo teria --- essas faculdades), mas emquanto é corpo \emph{vivo} e \emph{humano}, e O nosso corpo é \emph{vivo} e é \emph{humano} pela alma. Logo a alma é o \emph{principio} de \emph{todas} as faculdades humanas \footnote{A conclusão é um simples corollario das doutrinas expostas. A. alma é a forma substancial do corpo. Ora a forma substancial é o principio de todas as operações, e por isso de todas as energias ou faculdades do composto. Logo a alma é o principio de todas as faculdades humanas (\emph{Sum}. \emph{Theol}. p. L, q. 77, a. 5 ad 1). --- As faculdades não derivam da alma por uma mudança da propria alma, nem por uma sua acção physica e immeata, mas por um natural resultado, por uma necessaria emanação, assim como do sol emanam muitos raios. (Cf. \emph{Sum}, \emph{Theol}., l. c. a. 6). --- Repetimos \textbar{} que esta multiplice aptidão da alma não é contraria á sua \emph{simplicidade}.}. a 67. A alma não a o sujeito de todas as faculdades humanas. --- \emph{Sujeito} de uma faculdade é o ente, em que reside a propria faculdade. Ora a alma não é o ente, em que residem \emph{todas} as faculdades humanas. Com effeito, as. faculdades humanas são \emph{inorganicas} ou \emph{organicas}. As \emph{inorganicas} residem na alma, como em seu \emph{sujeito}; porque uma qualidade \emph{espiritual}, como é toda a faculdade \emph{inorganica}, só póde residir num sujeito \emph{espiritual}, como é a nossa alma. Às \emph{organicas} residem no \emph{composto} \pg{102} , como em seu \emph{sujeito}; pois que uma qualidade \emph{material}, como é toda a faculdade \emph{organica}, só pode residir num sujeito \emph{material}, como é o \emph{composto} de alma e de corpo. Logo a alma não é o \emph{sujeito} de todas as nossas faculdades \footnote{Diz S. Thomaz (\emph{Sum}. \emph{Theol}. p. 1, q. 77, a. 5): « Quaedam operationes sunt animae, quae exercentur per organa corporalia, sicut visio per oculum, auditus per aurem et ita de omnibus aliis operationibus nutritivae et sensitivae partis, et ideo potentiae, quae sunt talium operationum prineipia, sunt in \emph{conjuncto} sicut in \emph{subjecto}, e\& non in anima sola ».\par \emph{Spencer}, \emph{Taine,com} todos os \emph{positivistas} e \emph{materialistas}, pretendem que não só as faculdades vegetativas e as sensitivas mas tambem a intelligencia e a vontade consistem na conveniente disposição dos orgãos, e que por isso todas as faculdades humanas residem no \emph{composto}, como no seu proprio \emph{sujeito}. --- \emph{Platão}, \emph{Descartes}, \emph{Rosmini}, dizem que as faculdades \emph{sensitivas} residem na \emph{alma}, como no seu \emph{sujeito}.}.

68. As faculdades distinguem-se realmente da essencia da \textbf{alma}.

a) Um ente, que é \emph{acto}, distingue-se realmente do ente, que é \emph{potencia}; visto que uma coisa não póde ser, ao mesmo tempo e sob o mesmo aspecto, \emph{acto} e \emph{potencia}. Ora a essencia da alma é \emph{acto}, porque sempre \emph{actua}, isto é, aperfeiçõa e vivifica a materia, collocando-a numa determinada especie, ao passo que a faculdade é \emph{potencia}, porque não produz sempre a operação. Logo as faculdades distinguem-se \emph{realmente} da essencia da alma \footnote{Este argumento póde ser exposto de outro modo. --- Se a essencia da alma não se distinguisse realmente das suas faculdades, mas fosse ella mesma o principio \emph{immediato} das suas operações, seguir-se-hia que o ente dotado de alma, assim como é sempre vivo, assim tambem deveria produzir sempre as operações vitaes de que é capaz; como tambem, sempre e em cada momento, deveria sentir, entender e querer tudo o que póde sentir, entender e querer. Ora a experiencia attesta que a alma póde operar e deixar de operar, porque póde sentir e deixar de sentir, entender e deixar de entender; e as coisas, que ella sente, entende e quer, não as sente, entende e quer todas ao mesmo tempo. Logo as faculdades distinguem-se \emph{realmente} da essencia da alma (\emph{Cf}, \emph{Sum}. \emph{Theol}. p. T, q. 1591,}.

b) A substancia e o accidente são coisas \emph{realmente} distinctas entre si; porque a substancia não carece, como o accidente, de um sujeito a que adhira. Ora a alma é \emph{substancia}, e a faculdade é \emph{accidente}. --- A alma é \emph{substancia}, porque, como demonstrámos, a alma, sendo o principio \emph{primeiro} da vida, deve ser substancia. --- A faculdade é \emph{accidente}. Com effeito, a faculdade refere-se essencialmente Á operação como ao seu complemento, como a \emph{potencia} ao \emph{acto}. Ora a \emph{potencia} e o \emph{acto} pertencem ao \emph{mesmo genero}. Logo a faculdade e a operação \pg{103}  pertencem ao mesmo genero. Mas a operação pertence ao genero de \emph{accidente}; pois é coisa variavel e passageira. Logo tambem a faculdade pertence ao genero de \emph{accidente}. Logo as faculdades distinguem-se \emph{realmente} da essencia da alma \footnote{Neste argumento allegamos um axioma, que deve ser percebido e avaliado no seu verdadeiro sentido.\par O axioma é o seguinte: \emph{O acto e a potencia pertencem ao mesmo genero}. Isto quer dizer que, se a \emph{potencia} pertence ao genero de \emph{substancia}, tambem o \emph{acto} é \emph{substancia}, e vice-versa; e, se a \emph{potencia} pertence ao genero de \emph{accidente}, tambem o \emph{acto} deve ser \emph{accidente}, e vice-versa. Por quanto, o \emph{acto} determina e especifica a \emph{potencia}. Ora uma potencia \emph{substancial} não póde ser completada na sua \emph{especie} por um acto \emph{accidental}, pois um acto \emph{accidental} não completa a \emph{especie}, mas accrescenta-se a uma substancia jà completa na sua especie; --- nem uma potencia \emph{accidental} póde ser determinada e completada por um acto \emph{substancial}, porque um \emph{accidente} não póde sustentar uma \emph{substancia}, embora \emph{incompleta}. Logo, se a \emph{potencia é substancia}, tambem o seu \emph{acto é substancia}, e vice-versa; e, se a \emph{potencia} é \emph{accidente}, tambem o \emph{acto é accidente}, e vice-versa.\par Note se, porém, que este axioma --- a \emph{potencia} e \emph{o acto pertencem} ao mesmo \emph{genero} --- se refere \emph{exclusivamente} á potencia, que tem uma \emph{ordem essencial} e \emph{natural} ao seu acto, como ao seu \emph{fim adequado} e \emph{especificativo}, de modo que da potencia e do acto resulte um todo, composto de duas realidades, que se reduzem ao mesmo genero; como se verifica no caso presente, com relação ás faculdades ou potencias da alma. Por quanto, se a potencia não tem essa \emph{ordem natural} e \emph{essencial} ao acto, não se pode affirmar que acto e potencia pertencem ao mesmo genero; assim, quando se diz que a parede (\emph{potencia}) é alva (\emph{acto}), não se pode dizer que a parede e a alvura pertencem ao mesmo genero, e não se pode dizer porque esta \emph{substancia}, que é a parede, não se refere \emph{essencial} e \emph{naturalmente} a este \emph{accidente}, que é a alvura, nem é por este especificada.\par A. mesma conclusão, --- que em nós a faculdade, por ser uma coisa \emph{accidental}, se distingue.realmente da essencia da alma, --- dedul-a o Doutor Angelico de principios muito elevados e fecundos. Nas creaturas, diz S. Thomaz, o ser e a \emph{operação} são coisas realmente distinctas; porque só em Deus, cuja simplicidade e perfeição é infinita, o \emph{ser} e a \emph{operação} são uma e a mesma coisa. Ora o \emph{ser} é a actualidade da \emph{essencia}, e a \emph{operação} é a actualidade da \emph{potencia} ou \emph{faculdade}. Logo se, nas creaturas, o \emph{ser} e a \emph{operação} são coisas \emph{realmente} distinctas, tambem realmente distinctas devem ser a \emph{essencia} e a \emph{potencia} ou \emph{faculdade}. Por quanto diversas actualidades, ou diversos actos, suppõem e exigem diversos sujeitos, pois deve haver proporção entre o acto e o proprio sujeito. (\emph{Sum}. \emph{Theol}., p. 1, q. 94, a. 1 e 2; q. '(9, a. 1; \emph{De spirit}. \emph{ereat}., a. 11; \emph{De anima}, a. 12).\par D'onde se vé a razão pela qual os \emph{pantheistas} impugnam a nossa these. Na verdade, se se admitte que nas creaturas o \emph{ser} e a \emph{operação} são uma e a mesma coisa, attribue-se á creatura o que é exclusivamente proprio de Deus, e não é diffieil concluir, como elles concluem, que a nossa alma não se distingue realmente de Deus.}.

o) Duas coisas distinguem-se \emph{realmente} entre si, quando têm \pg{104}  notas ou caracteres oppostos. Ora a \emph{faculdade} e a \emph{essencia} da alma têm notas ou caracteres oppostos. Com effeito, a faculdade é multipla, relativa, incompleta, passa da potencia para o acto, e é capaz de maior ou menor desenvolvimento, segundo a diversa indole do sujeito em que reside, --- ao passo que a essencia da alma é unica, absoluta, completa, não admitte mudança, nem é capaz de progresso, mas é sempre uma e a mesma. Logo as faculdades distinguem-se \emph{realmente} da essencia da alma. --- Se as faculdades emanam necessariamente da alma, mas não são a sua essencia, deve dizer-se que são suas \emph{propriedades} \footnote{Todos os philosophos concordam em que existe uma certa distincção entre a essencia da alma e as suas faculdades, mas discordam quando se trata de determinar a natureza d'essa distineção.\par Uns, como \emph{Aristoteles}, \emph{Santo Thomas}, \emph{Leibnitz}, e os \emph{Escholasticos}, afirmam que essa distincgão é \emph{real}, isto é, que a essencia da alma e as suas faculdades são varias entidades distinctas; por quanto as faculdades, derivando da essencia, são o principio \emph{procimo} das operações, e a essencia é o seu principio \emph{remoto}. --- E' a doutrina que seguimos e defendemos.\par Outros, os \emph{Cartesianos}, seguindo os \emph{Pantheistas}, dizem que a essencia da alma e as suas faculdades não são entidades \emph{realmente} distinctas, mas que são identicas, pois a essencia da alma é o principio \emph{provimo} das operações, e as faculdades são a mesma essencia, diversamente modificada.\par \emph{A} opinião dos \emph{Cartesianos} é falsa, não só pelas razões adduzidas no texto, mas tambem pelas seguintes:\par a) \emph{Basea-se num falso fundamento}. Por quanto, erradamente nega uma distincção \emph{real} entre a substancia e o accidente.\par \emph{b}) \emph{E}' \emph{contraria dá experiencia}. Na verdade, a consciencia attesta que as faculdades estão por tal forma relacionadas entre si, que uma depende de outra, uma regula e domina outra, uma está em lucta com a outra. Assim, o exercicio da vontade depende da intelligencia, esta depende dos sentidos, e os sentidos das faculdades vegetativas; a razão domina o appetite sensitivo, e a vontade regula todas as outras faculdades; a razão ás vezes está em lucta com a imaginação, e á vontade com as pelkóon Ora, se as faculdades não se distinguissem \emph{realmente} da essencia da alma, uma e a mesma simplicissima essencia dependeria de si, regularia e dominaria a si mesma, e estaria em lucta comsigo; o que é absurdo. Logo as faculdades distinguem-se \emph{realmente} da essencia da alma (\emph{Of}. S. Th. \emph{De Sp}. \emph{creat}. a. 11).\par \emph{c}) \emph{Leva a falsas consequencias}. Se a essencia da alma não se distinguisse \emph{realmente} das suas faculdades, --- ou deveriamos admittir uma unica faculdade, porque unica é a essencia, --- ou deveriamos dizer que a essencia da alma consiste na \emph{collecção} das varias faculdades. Ora ambas estas consequencias são falsas. --- E' falsa \emph{a primeira}. Na verdade, aquella unica faculdade ou seria \emph{organica}, ou \emph{inorganica}. Se fosse \emph{organica}, não poderia exercer as operações \emph{espirituaes}, como são as da intelligencia e da vontade. Se fosse \emph{inorganica}, não poderia exercer as operações \emph{organicas}, como são as vegetativas e as sensitivas. --- E' falsa a \emph{segunda}. Com effeito, nessa hypothese\pgn{105}, a alma não seria \emph{substancia}, mas \emph{accidente}, porque uma \emph{collecção} de \emph{accidentes}, como são as faculdades, não pode deixar de ser \emph{accidente}.\par Objectam: Se as faculdades se distinguissem \emph{realmente} da essencia da alma, poderiam separar-se d' ella. Ora não podem separar-se. Logo.\par A objecção parte do supposto de que um ente, que se distingue \emph{realmente} de outro, deva poder separar-se d' este. Ora tal supposto póde ser verdadeiro, mas póde tambem ser falso. --- E' verdadeiro, quando se trata de duas coisas, que não só são \emph{realmente} distinctas entre si, mas que tambem podem ambas continuar a subsistir quando separadas uma de outra. --- E' falso, quando se trata de duas coisas, que, se são \emph{realmente} distinctas entre si, não são ambas \emph{subsistentes}; assim o esplendor do sol e o proprio sol são duas coisas \emph{realmente} distinctas, e todavia o esplendor não póde separar-se do sol, porque o esplendor não é \emph{subsistente}. E' o que acontece no nosso caso. As faculdades são \emph{realmente} distinctas da essencia da alma, mas não podem separar-se d'ella, porque são \emph{accidentes}, e o accidente exige naturalmente um sujeito a que adhira.\par Insistem: As faculdades, por isso mesmo que são \emph{accidentes}, podem separar-se da essencia da alma, pois o \emph{accidente} póde separar-se da substancia.\par Ha duas especies de \emph{accidente}: \emph{lógico} e \emph{ontologico}. --- Accidente \emph{logico} é o que não deriva \emph{necessariamente} da essencia, mas póde indifferentemente dar-se ou deixar de dar-se. Este accidente póde, ás vezes, separar-se do sujeito a que adhere. --- Accidente \emph{ontologico} é o ente que precisa de sujeito a que - ndhira, quer emane \emph{necessariamente} da essencia do proprio sujeito, quer não. Este accidente, --- quando emana \emph{necessariamente} da essencia, diz-se \emph{propriedade} e não póde separar-se naturalmente da mesma essencia, --- mas, quando emana \emph{contingentemente} da essencia, póde separar-se d' ella. Ora á faculdade não é \emph{accidente logico}, mas \emph{ontologico}, e emana \emph{necessariamente} da essencia, é uma \emph{propriedade}, e não póde separar-se naturalmente da essencia mas onde se encontra esta, deve encontrar-se aquella (S. Th., l. e. ad 1).\par Replicam: Não póde conceber-se a essencia sem as suas faculdades. Logo aquella não se distingue \emph{realmente} d' estas.\par Respondemos que, embora a alma não possa existir sem as suas faculdades, especialmente sem a intelligencia e á vontade, todavia podemos conceber o que é a alma sem faculdades; porque podemos concebel-a como um sujeito, constituido numa certa e determinada especie, existente em si, principio e raiz de que naturalmente derivam todas as faculdades. Na verdade, primeiramente forma-se a idea da essencia de uma coisa, e depois a idea dos seus attributos ou propriedades.}.

%%>>> art-14

\section*{Artigo III.}
\addcontentsline{toc}{section}{Artigo III.}

\pg{105} \textbf{As faculdades consideradas em relação}

\textbf{ás operações}.

%%>>> t69-2

\tese{69}{As faculdades distinguem-se realmente das suas operações.}

a) A causa distingue-se \emph{realmente} do seu effeito. Ora as faculdades são as causas, e as operações são seus effeitos. Logo as faculdades distinguem-se \emph{realmente} das \emph{suas operações}.

b) A consciencia attesta que as operações mudam constantemente, se suecedem umas ás outras, e todavia a faculdade, que as produz, fica sempre a mesma. Ora isto seria impossivel, se as faculdades não fossem \emph{realmente} distinctas das suas operações. Logo \footnote{\emph{Descartes}, assim como tinha negado uma distincção \emph{real} entre a essencia da alma e as suas faculdades, assim tambem nega uma distincção \emph{real} entre as faculdades e as suas operações, e diz que a faculdade está sempre em exercicio, porque uma faculdade, que não fosse sempre activa, seria uma contradicção.\par A opinião de \emph{Descartes} é falsa. O que é essencial á potencia, não é o exercicio \emph{actual} da sua força, mas é o \emph{poder} de operar, --- poder, que a faculdade habitualmente possue, mesmo quando não opéra.\par Os \emph{positivistas} negam a distineção \emph{real} entre as faculdades e as operações, e sustentam que as faculdades são apenas palavras ou conceitos, de que nos servimos para agrupar e significar uma certa serie de phenomenos e de operações. Assim, dizem elles, os actos da intelligencia, por causa da sua similhança, são denotados pelo mesmo nome e attribuidos a0 mesmo . principio \emph{hypothetico} --- á intelligencia; embora a intelligencia seja apenas a serie dos actos intellectuaes. Diga-se o mesmo em relação á vontade.\par Como so vê, a opinião dos \emph{positivistas} é contraria á razão e á experiencia.\par Em \textsc{Deus}, pela sua infinita perfeição e simplicidade, as operações não se distinguem realmente das faculdades, nem estas da essencia.}.

%%>>> t70

\pg{106} \tese{70}{As faculdades têm uma tendencia natural para as suas operações.}
Todo o ente tem uma tendencia natural para o seu fim; pois o fim é um bem, que aperfeiçoa o ente que o consegue. Ora o fim da faculdade é a operação. Logo as faculdades têm uma tendencia natural para as suas operações. Assim os olhos tendem naturalmente para as \emph{côres}, o ouvido para os \emph{sons}, o gosto para os \emph{sabores}, ete. \footnote{Cf. S. Th. \emph{De verit}., q. 22, a. 1. --- Esta \emph{tendencia} das faculdades chama-se tambem \emph{appetite natural} ou \emph{intenção}, porque é em virtude della que as faculdades se inclinam naturalmente para o seu /\emph{im}, que consiste no exercicio das proprias operações.\par Os principaes effeitos d'esta \emph{tendencia natural} são dois: --- 1.9) Quanto mais intensa fór a tendencia da faculdade, tanto mais perfeito será o acto, porque o effeito deve ser proporcionado á causa; assim o acto intellectual será tanto mais perfeito, quanto maior fôr a attencaão da intelligencia. --- 2.º) Quanto mais intenso fôr o acto de uma faculdade, tanto mais remisso será o acto das outras faculdades, pois que a força da alma é finita, e por isso, se ella augmenta numa faculdade, deve diminuir noutra; assim, se a \emph{vista} se fixar muito intensamente num objecto, o \emph{ouvido} não perceberá o que se está dizendo (Cf. S. Th. \emph{De} Ver.. q. 18, a. 3).}.

%%>>> t71

\tese{71}{As faculdades e o habito.}
As faculdades não só têm uma \emph{inclinação} natural para as respectivas operações, mas \pg{107}  podem ser dotadas de uma certa \emph{disposição}, que torne mais facil o seu exercicio. Esta \emph{disposição} chama-se \emph{habito} \footnote{O \emph{habito} distingue-se --- não só da \emph{tendencia}, de que as faculdades são naturalmente dotadas, --- mas tambem das proprias \emph{faculdades}. --- Distingue-se da \emph{tendencia}. Na verdade, a \emph{tendencia} deriva da natureza, convém a \emph{todas} as faculdades da alma e a todos os entes, e inclina a faculdade a produzir \emph{todos} os actos que lhe são proprios; ao passo que o \emph{habito} póde derivar e muitas vezes deriva da repetição dos actos, não convém a todas as faculdades nem a todos os entes, e inclina a faculdade a operar mais de um modo do que de outro. --- Distingue-se da \emph{faculdade}. Com effeito, a faculdade deriva da essencia do ente, é por si indifferente para produzir este ou aquelle acto, não póde separar-se da essencia de que deriva; ao passo que o \emph{habito} deriva ou immediatamente de \textsc{Deus} ou da repetição dos actos, dispõe a faculdade para operar mais de um modo do que de um outro, e póde separar-se da faculdade (\emph{Cf}. \emph{Sum}. \emph{t}/., 1. 2., q. 49-54).}.

Digamos alguma coisa ácerca do \emph{habito}, seu \emph{sujeito} e \emph{divisão}.

a) Habito. --- \emph{Habito} é uma qualidade estavel ou permanente, que dispõe a faculdade a operar mais de um modo do que de outro, bem ou mal, conforme a operação convém ou repugna á natureza do agente. E' o ultimo complemento da faculdade \footnote{\emph{Q habito} divide-se em varias especies, conforme se considera no seu \emph{sujeito}, na sua \emph{origem} ou nos seus \emph{actos}.\par \emph{a}) Considerado no seu \emph{sujeito}, o \emph{habito} divide-se em \emph{intellectual} e \emph{moral}. --- \emph{Intellectual} é o que reside na \emph{intelligencia}, como o habito da \emph{sciencia}. --- \emph{Moral} é o que reside na \emph{vontade}, como o habito da \emph{virtude}.\par \emph{b}) Considerado na sua \emph{origem}, o \emph{habito é natural} ou \emph{adquirido}. --- \emph{Natural} é o que deriva de uma inclinação natural, sem esforço das faculdades intellectuaes; tal é o habito que nos dispõe a perceber prompta e facilmente os \emph{primeiros principios} da razão. --- \emph{Adquirido} é o que provém da repetição do acto; como é o habito da \emph{temperança}.\par \emph{e}) Considerado'no seu \emph{acto}, o \emph{habito} poder ser \emph{natural} ou \emph{sobrenatural}. --- E \emph{natural}, se dispõe para a producção de um acto, que não excede as forças proprias dos agentes creados; tal é o habito da \emph{sciencia}. --- E' \emph{sobrenatural}, se dispõe para a producção de um acto, que excede as forças proprias dos agentes creados; tal é o \emph{habito} da \emph{F}6, da \emph{Esperança}, e da \emph{Caridade}. --- O habito \emph{sobrenatural} chama-se tambem \emph{infuso}, porque é \textsc{Deus} quem immediata e directamente o infunde na alma. Mas, então, considera-se na sua \emph{origem},}.

b) Sujeito do habito. --- Sujeito do \emph{habito é exclusivamente} a \emph{intelligencia} e á \emph{vontade}. Com effeito, estas duas faculdades, por serem indifferentes em relação ás suas operações e aos seus objectos, seriam indeterminadas e muito vagarosas na sua actividade, se não fossem dispostas pelo \emph{Habito} a operar mais de um modo, do que de outro. --- Além d isso, estas duas faculdades podem adquirir por si mesmas o \emph{habito}; porque, sendo dotadas de reflexão e de liberdade, podem mudar e regular a direcção da sua actividade.

Dissemos --- \emph{exclusivamente}; pois as faculdades \emph{vegetativas} e as \emph{sensitivas} não \emph{precisam} de \emph{habitos}, nem os \emph{podem} adquirir. - Não \emph{precisam} de \emph{habitos}; porque, pela \emph{natureza}, são determinadas a opezar sempre do mesmo modo. - Nào os \emph{podem} adquirir; porque, não tendo reflexão nem liberdade, não podem regular a sua actividade \footnote{\emph{Garnier}, \emph{Maine de Biran} e outros escriptores modernos attribuem o \emph{habito} não só á intelligencia e á vontade, mas tambem ás faculdades \emph{sensitivas}, ao corpo humano, e a todos os seres organisados.\par Respondemos que o \emph{Labito}, pelas razões adduzidas no texto, compete \emph{exclusivamente} á intelligencia e á vontade. --- Todavia concordamos em que, num \emph{sentido mais lato}, o \emph{habito} póde attribuir-se ás nossas faculdades \emph{sensitivas}, \emph{e} aos animaes. Mas, então, as nossas faculdades \emph{sensitivas} e os animaes só podem chamar-se \emph{sujeitos} do \emph{habito}, emquanto são dirigidos pela razão \emph{humana} a operar sempre ou frequentemente de um certo modo e não de outro; por fórma que o /\emph{iabito} é sempre um effeito da \emph{razão}.}.

c) Effeitos do habito. --- Pelo \emph{Habito} o exercicio das faculdades torna-se \emph{uniforme}, \emph{agradavel} e \emph{facil}. --- \emph{Uniforme}; porque o \emph{habito}, por ser uma qualidade \emph{permanente}, dispoe a operar mais de um modo que de um outro, e por isso chama-se uma \emph{segunda natureza}. --- \emph{Agradavel}; porque o \emph{habito} torna a operação quasi \emph{natural}, isto é, conveniente á natureza, e tudo o que convem á natureza é motivo de agrado. --- \emph{Facil}; porque o agente, sem muita difficuldade, executa uma operação, para a qual está bem disposto \footnote{Devemos advertir que o \emph{Habito}, embora torne mais \emph{facil} o acto, todavia não tira a liberdade de fazer o contrario, e que muitas vezes, longe de tornar o acto mais perfeito, prejudica-o, porque diminue a attencüo e a reflexão da intelligencia.}.

%%>>> cap-3

\chapter*{CAPITULO SEGUNDO}
\addcontentsline{toc}{chapter}{CAPITULO SEGUNDO}

\pg{108} \textbf{FACULDADES} SENSITIVAS DO HOMEM \textsc{SUMMARIO}: --- Sentidos externos. --- Sentidos internos. --- Appetite sensitivo,

faculdade locomotiva e vocal.

%%>>> art-15

\section*{Artigo I.\\ Sentidos externos.}
\addcontentsline{toc}{section}{Artigo I. --- Sentidos externos.}

%%>>> t72

\tese{72}{Faculdades sensitivas do homem.}
Faculdades \emph{sensitivas} do homem são aquellas por meio das quaes elle attinge objectos \emph{materiaes} na sua \emph{singularidade}, isto é, cercados de condições \pg{109}  \emph{individuaes}. --- O homem attinge tambem objectos \emph{immateriaes}, e objectos \emph{materiaes} na sua \emph{universalidade}, mas só por meio das suas faculdades \emph{intellectuaes} \footnote{Não tratamos aqui de proposito das faculdades \emph{vegetativas} do homem, não só porque no artigo Il do capitulo I da Secção primeira d'este tratado fallámos Acerca dos orgãos da vida \emph{vegetativa} e das suas respectivas funcções, mas tambem porque o argumento não é tão importante que exija um maior desenvolvimento. Tratamos, porém, das faculdades \emph{sensitivas}, não só porque as suas operações estão sujeitas ao alcance da consciencia e assim podem mais facilmente ser percebidas e analysadas, mas tambem, e sobretudo, porque as faculdades \emph{sensitivas} estão ligadas com as \emph{intellectuaes}.}.

%%>>> t73

\tese{73}{Todas as faculdades sensitivas são materiaes.}

a) Uma faculdade é \emph{material}, quando o seu objecto é \emph{material}; pois, como provámos, as faculdades conhecem-se e especificam-se pelos seus \emph{objectos}. Ora o objecto de todas as faculdades \emph{sensitivas é material}; porque ou é materia, ou é cercado de condições materiaes. Logo todas as faculdades \emph{sensitivas} são \emph{materiaes}.

b) Uma faculdade diz-se \emph{material}, quando o seu sujeito é \emph{material}, isto é, composto de materia e de fórma. Ora o sujeito de todas as faculdades \emph{sensitivas} --- não é só o \emph{corpo}, porque o corpo, como tal, não é capaz de sentir (aliás todo o corpo seria dotado de sensibilidade), --- nem é só a \emph{alma}, porque a nossa alma, sendo \emph{espiritual} ou \emph{immaterial}, só por si não pode perceber immediata e directamente objectos \emph{materiaes} na sua \emph{singularidade} (visto que objectos materiaes e singulares não podem actuar numa faculdade immaterial), --- mas é o \emph{composto} de alma e de corpo, isto é, é o \emph{orgão amimado}. Logo todas as faculdades \emph{sensitivas} são \emph{materiaes} \footnote{Um corpo não pode actuar ou fazer impressão no espirito; porque, sendo \emph{extenso}, o corpo opera por movimento local, e por isso só pode actuar num sujeito tambem extenso. E' o adagio repetido muitas vezes por S. Thomaz: \emph{Nihil corporeum imprimere potest in rem incorpoream} (\emph{Sum}. \emph{Theol}. p. L, q. 84, a. 6) e por isso \emph{omnis operatio sensitivae partis est conjuncti} (p. T, q. (b, á. 5) --- Alem d'isso, se as faculdades \emph{sensitivas} fossem \emph{immateriaes}, os Anjos seriam dotados tambem de vida \emph{sensitiva}, e as almas dos brutos seriam \emph{espirituaes} e \emph{immortaes}: o que se não pode admittir.\par Objectam: A faculdade sensitiva é simples e inextensa. Ora uma propriedade simples e inextensa não pode ser recebida no orgão. Logo a faculdade sensitiva não é organica, não é material.\par Resposta. Toda a faculdade sensitiva é por si simples e inextensa, porque deriva da alma, e porque é forma do orgão, e toda a forma, substancial ou accidental, é \emph{por si} simples e inextensa (só o corpo é por si extenso, isto é, dotado de quantitade). Ora a faculdade sensitiva, embora por si seja\pgn{110}simples e inextensa, pode comtudo ser recebida num sujeito extenso, como toda e qualquer forma, e então torna-se extensa não \emph{por si}, mas \emph{por causa do sujeito} ou \emph{orgão}, que ella informa. Mas, sendo forma de um orgão diversamente organisado, a faculdade sensitiva não se divide em partes, quando em partes se divide o proprio orgão; assim se dividirmos em duas partes o orgão da vista, não se divide a energia visual. --- Note-se que \emph{faculdade sensitiva} pode significar e significa --- tanto a energia sensitiva, --- como o orgão, isto é, a parte especial do corpo, que a energia informa, ---- e tambem. o \emph{composto} da energia e do orgão.}.

%%>>> t74

\pg{110} \tese{74}{Especies das faculdades sensitivas do homem.}
Às especies das faculdades \emph{sensitivas} do homem são \emph{quatro}, a saber: --- a \emph{perceptiva}, que abrange os \emph{sentidos externos} e os \emph{internos}, --- a \emph{appetitiva}, --- a \emph{locomotiva}, --- e a \emph{vocal}. Com effeito, o homem, por meio da \emph{sensibilidade}, ou \emph{percebe} um objecto \emph{material}, como \emph{tal} (fae. \emph{perceptiva}), --- ou \emph{tende} para elle por uma inclinação intrinseca (fac. \emph{appetitiva}, ou \emph{espontaneamente} se move do seu logar para o alcançar (fae. \emph{locomotiva}); ou manifesta, pela \emph{voz}, o que se passa no seu interior (fae. \emph{vocal}) \footnote{Estas especies de faculdades sensitivas estão por tal fórma Telacionadas entre si, que uma exige outra, uma não póde estar sem outra. Assim a faculdade \emph{perceptiva} exige a \emph{appetitiva}, porque a toda a percepção succede uma tendencia; e, vice-versa, a faculdade \emph{appetitiva} suppõe a \emph{perceptiva}, porque é impossivel tender para um objecto que não se conhece. Como tambem, a faculdade \emph{perceptiva} e á \emph{appetitiva} exigem a faculdade \emph{locomotiva}, porque seria inutil perceber e desejar um bem, se não houvesse meio para o alcançar; e, vice-versa, a faculdade \emph{locomotiva} suppõe a faculdade \emph{perceptiva} e á \emph{appetitiva}, porque o homem não poderia mover-se de um para outro logar para alcançar um bem, se este bem ríáo fosse previamente conhecido e desejado. Diga-se \emph{o} mesmo com relação á faculdade \emph{vocal}. :}.

Neste artigo tratamos dos \emph{sentidos externos}, --- no seguinte, dos \emph{sentidos internos}, --- e no terceiro, do \emph{appetite sensitivo} e da faculdade \emph{locomotiva} e \emph{vocal}.

%%>>> t75

\tese{75}{Sentidos externos.}
Sentidos \emph{externos} são os que existem na superficie \emph{externa} do organismo humano, e por cujo intermedio o homem percebe os seres \emph{materiaes} e \emph{externos}. --- Pelos sentidos \emph{externos} percebemos os seres \emph{materiaes} e \emph{externos} por causa das propriedades \emph{sensiveis}, de que estes são dotados e que são o objecto \emph{proprio} e \emph{directo} dos mesmos sentidos \footnote{O objecto dos sentidos, como já dissemos (na \emph{Cosmologia especial}), é triplice: \emph{proprio}, \emph{commum}, e \emph{concomitante}. --- O \emph{proprio} é o que póde ser percebido por um só sentido, e que \emph{por si} e \emph{immediatamente} impressiona o orgão: tal é toda a \emph{qualidade sensivel} dos corpos, assim a cór é o objecto \emph{proprio} dos \emph{olhos}, \emph{o som} do \emph{ouvido}, ete. --- O commum é o que é percebido por mais de um sentido; tal é a \emph{figura}, \emph{á grandeza}, o \emph{movimento}, etc. --- O \emph{concomitante} é o que não é percebido por si mesmo, mas por meio de\pgn{111}outro objecto, a que está unido; tal é a \emph{substancia corporea}, que não é percebida em si mesma, mas nas \emph{qualidades sensiveis}, que a envolvem.\par \emph{Reid}, \emph{Rosmini} e \emph{Gioberti} sustentam que a \emph{substancia corporea} não póde de \emph{modo algum} ser percebida pelos sentidos, mas só pela intelligencia.\par \emph{A} opinião d'esses philosophos é inacceitavel. Na verdade, por meio das faculdades sensitivas percebemos as qualidades corporeas, não \emph{éndeterminadas} e \emph{abstractas}, mas \emph{determinadas} e \emph{concretas}; assim os olhos não vêm a cór, mas \emph{esta cór}. Ora as qualidades são \emph{determinadas} e \emph{concretas} emquanto adherem a um determinado sujeito ou a uma determinada substancia. Logo, por meio das faculdades sensitivas, \emph{directamente} percebemos as qualidades corporeas, --- mas, \emph{indirectamente} ou por \emph{concomitancia}, embora de \emph{um modo} confuso, attingimos a propria substancia, a que as qualidades adherem.\par Nem se diga que só pela intelligencia podemos perceber a \emph{substancia}. Pela intelligencia percebemos a \emph{essencia} da substancia, mas pelos sentidos percebemos a substancia como uma coisa \emph{concreta} e \emph{singular}, da qual ignoramos a \emph{essencia}, mas que é o sujeito das qualidades sensiveis.}.

%%>>> t76

\pg{111} \tese{76}{Numero dos sentidos externos.}
Os sentidos \emph{externos} do homem são \emph{cinco}: a \emph{vista}, 0 \emph{ouvido}, o \emph{olfacto}, e \emph{gosto}, e o \emph{tacto}.

a) Vista. --- \emph{Vista} é o sentido pelo qual o homem percebe as \emph{córes}. --- Os orgãos da \emph{vista} são os \emph{olhos}.

b) Ouvido. --- \emph{Ouvido} é o sentido pelo qual \emph{o} homem percebe os \emph{sons}, produzidos pelos movimentos vibratorios dos corpos elasticos. --- O orgão d'este sentido é o \emph{ouvido}.

c) Olfacto. --- \emph{Olfacto} é o sentido pelo qual o homem percebe os \emph{cheiros}. --- A sua séde é principalmente a \emph{pituitaria}, que guarneee as fossas \emph{wasaes}.

d) Gosto. --- \emph{Gosto} é o sentido pelo qual o homem percebe os \emph{sabores}. --- O seu orgão principal é a \emph{lingua}.

e) Tacto. --- \emph{Tacto} é o sentido pelo qual o homem percebe os objectos, que pela sua resistencia fazem nelle impressão. --- A sede principal do \emph{tacto} encontra-se na superficie palmar das extremidades dos dedos, na parte anterior da bocca e na ponta da lingua \footnote{Não damos ao assumpto um desenvolvimento maior; pois que, no artigo III do Capitulo I da \emph{primeira} secção, apresentámos as principaes noções anatomicas e physiologicas ácerca dos \emph{cinco} sentidos \emph{externos}.\par Só queremos accrescentar alguma coisa ao que dissemos com relação á \emph{vista}, Quem considera as partes \emph{essenciaes} e \emph{accessorias} dos \emph{olhos}, vé logo que cada uma d'essas partes tem o seu fim parcial, e o orgão todo tem um fim completo, que é a perfeita percepção do seu objecto. \emph{Tyndall} diz: « Como instrumento pratico, e attendendo ás adaptações que neutralizam os seus defeitos, o olho será sempre uma \emph{maravilha} para toda a mente capaz de reflexão » (\emph{Da lus}) Esta ordem maravilhosa, que a estructura do olho revela, não pode ter outra causa a não ser uma intelligencia, e uma intelligencia infinita, dotada de poder infinito. Por isso o grande anatomista\pgn{112}\textbar{} \emph{Falloppio}, que foi o primeiro a estudar os musculos do olho humano, conclue a sua obra \emph{De musculis oculorum} por estas palavras: « Junctis ad sidera palmis, Deo optimo et maximo, qui in homine oculos tam insigni magisterio exornavit, hymnum laudis canto ».}.

%%>>> t77

\pg{112} \tese{77}{Os sentidos externos não são mais que os cinco enumerados.}

a) Os sentidos \emph{externos} do homem serão tantos, quantas forem as diversas especies dos objectos, que por elles pode- \textbar{} mos perceber; pois as faculdades distinguem-se, ou especifi- \textbar{} eam-se, em ultima analyse, pela diversidade dos seus \emph{objectos} \textbar{} \emph{formaes}. Ora as diversas especies dos objectos, que pelos sentidos \emph{externos} podemos perceber, são cinco: as \emph{córes}, OS \emph{SONS}, os \emph{cheiros}, os \emph{sabores}, e as qualidades relativas ao \emph{facto}; pois tudo o que póde impressionar os sentidos \emph{externos} reduz-se a \emph{k} uma d' estas \emph{cinco} especies de objectos. Logo os sentidos \emph{exteri nos} não são mais que os \emph{cinco} enumerados. \textbar{}

b) Os sentidos \emph{externos} do homem se vão tantos, quantos forem os orgãos \emph{externos}, aptos para a sensação; pois a cada \textbar{} orgão, apto para a sensação, corresponde um sentido distincto e especial. Ora os orgãos \emph{externos}, aptos para a sensação, são \emph{cinco}, a saber: o \emph{visual}, O \emph{auditivo}, o \emph{olfactivo}, o \emph{gustativo}, e O orgão do \emph{tacto}. Logo os sentidos \emph{erternos} do homem não são mais que os \emph{cinco} enumerados \footnote{Os homens estão convencidos de que os sentidos \emph{externos} não são mais que cinco: \& \emph{vista}, o \emph{ouvido}, o \emph{olfacto}, o \emph{gosto} e o \emph{tacto}.\par Todavia ha excepções nesta convicção.\par a) \emph{Locke} (\emph{De int}. \emph{hum}., lib. II, c. 2, 88) duvidou de que no homem nao houvesse um outro sentido, distincto e diverso dos cinco enumerados.\par \emph{b}) Os sequazes do \emph{magnetismo animal} ou \emph{Aypnotismo} sustentam que a alma, durante o somno magnetico, desenvolve algumas faculdades sensitivas externas, que não são reduotiveis a nenhum dos cinco sentidos.\par \emph{c}) Finalmente \emph{Balmes} (\emph{Filos}. \emph{fund}., t. I, lib. 2, c. 17) ensina que o homem, com o andar do tempo, póde adquirir um novo sentido externo.\par Estas opiniões são falsas.\par \emph{a}) E' falsa a opinião de \emph{Locke}. Com effeito, para que podesse existir um outro sentido \emph{externo}, distincto e diverso dos cinco enumerados, seria necessario quo nos corpos existisse uma qualidade sensivel, distincta e diversa da \emph{cór}, dos \emph{sons}, dos \emph{cheiros}, dos \emph{sabores}, \emph{e} das qualidades relativas ao \emph{tacto}; pois que \emph{as} faculdades distinguem-se pelos seus objectos. Ora, como o proprio \emph{Locke} confessa, nos corpos não existe nenhuma qualidade sensivel, que não seja reductivel a uma ou outra das cinco enumeradas. Logo é impossivel que exista um sentido externo, distincto e diverso dos cinco enumerados.\par \emph{b}) E' falsa a opinião dos sequazes do \emph{magnetismo animal}. Por quanto, as faculdades sensitivas externas, que o homem desenvolveria durante o somno magnetico, seriam \emph{naturaes}, e todavia só produziriam o seu acto\pgn{113}num estado \emph{anormal}. Ora isto é absurdo; porque o que é natural deve revelar-se no estado \emph{regular}, e não num estado \emph{anormal}, --- Accrescente-se que o somno \emph{magnetico}, como todo o estado insolito da alma, se póde augmentar ou diminuir a energia de uma faculdade, não póde produzir a mesma faculdade.\par \emph{c}) E' falsa a opinião de \emph{Balmes}. Na verdade, se o homem, com o andar do tempo, podesse adquirir um novo sentido externo, \& sua essencia estaria sujeita á mudança; pois que as faculdades derivam naturalmente da essencia, e novas faculdades suppõem nova essencia. Ora a essencia das coisas é immutavel. Logo é impossivel que o homem, com o andar do tempo, adquira um novo sentido externo. --- Além d' isso, se o homem podesse, com o andar do tempo, adquirir um novo sentido, deveriamos dizer que Deus não concedeu ao homem, desde o principio, \emph{todas} as propriedades \emph{aaturaes}, como são os sentidos. Ora Deus não nega ás coisas o que é proprio da natureza d'ellas. Diz S. Thomaz: « Deus, qui est institutor naturae, non subtrahit rebus, quod est proprium naturis earum » (\emph{Contr}. \emph{Gent}., l. II, c. 55; Cf. \emph{Sum}. \emph{Theol}., p. Ij q. \&8, a. 3).\par Devemos advertir que, embora o homem viesse a adquirir, com o andar do tempo, um novo sentido \emph{externo}, comtudo nem por isso tornar-se-hia incerto ou falso o testemunho dos nossos sentidos, como pretende \emph{Lamennais}. Os conhecimentos adquiridos pelo \emph{sexto} sentido e os adquiridos por meio dos outros \emph{cinco} sentidos seriam \emph{diversos}, mas não \emph{contrarios}; assim como os conhecimentos dos \emph{cinco} sentidos são \emph{diversos} entre si, mas não \emph{contrarios}.\par 8 Vor. II.}.

%%>>> t78

\pg{113} \tese{78}{Comparação dos sentidos externos.}
Os sentidos \emph{externos} podem considerar-se --- emquanto á sua \emph{excellencia}, --- emquanto á sua \emph{certeza}, --- e emquanto á sua \emph{vantagem}.

a) Emquanto á \emph{excellencia}, um sentido é mais nobre do que outro, quando o objecto d'aquelle é mais \emph{extenso} que o d'este. Por isso o mais nobre entre todos os sentidos é a \emph{vista}, pela qual \emph{directa} e \emph{immediatamente} percebemos as \emph{córes}, e \emph{indirecta} e \emph{mediatamente} attingimos a \emph{figura}, a \emph{grandeza}, o \emph{movimento}, \emph{a quietação}, ete. Menos nobre do que a \emph{vista} é O \emph{ouvido}, que se limita á percepção dos \emph{sons}. Em seguida vem o \emph{o/facto}, \emph{e} finalmente o \emph{gosto} e o \emph{tacto}.

b) Emquanto á \emph{certeza}, o \emph{tacto} e a \emph{vista} são superiores a todos os outros sentidos, e por isso, querendo exprimir o mais elevado grau de certeza, dizemos: \emph{íoca-se com}. \emph{as mãos}: \emph{metté-se pelos olhos dentro}. O testemunho do \emph{facto} é mais seguro que o da \emph{vista}: pois esta ás vezes está sujeita a illusões.

c) Emquanto á \emph{vantagem}, é necessario distinguir. Ou se trata da vantagem relativa á vida \emph{intellectual} do homem, ou da vantagem relativa á sua vida \emph{organica}. --- Se se trata da vantagem relativa á vida \emph{intellectual}, \emph{á vista} e o \emph{ouvido} são superiores aos outros sentidos: pois a \emph{vista} ajuda a conhecer as differenças das coisas, e o \emph{ouvido} serve para perceber o \pg{114}  ensino dos sabios. --- Se se trata da vantagem relativa á vida \emph{organica}, o \emph{tacto} e o \emph{gosto} são superiores aos outros sentidos \footnote{Dissemos que, se se trata da vantagem relativa á vida \emph{intellectual}, a \emph{vista} e o \emph{ouvido} são superiores aos outros sentidos. E, sempre sob este aspecto, o \emph{ouvido} é superior á \emph{vista}. Por quanto o homem é essencialmente \emph{discipulo}, o numero dos conhecimentos que recebemos dos outros é superior ao numero do conhecimentos que adquirimos pelos nossos esforços, e por isso o sentido do \emph{ouvido} é mais necessario para o desenvolvimento da intelligencia do que o sentido da \emph{vista}. (S. Thomaz, \emph{de sens}., lect. ID).\par Emquanto ao \emph{desenvolvimento}, primeiramente é o \emph{tacto}, que é o fundamento de todos os sentidos, depois o \emph{gosto}, em seguida o \emph{olfacto} e a \emph{vista}, e finalmente o \emph{ouvido}. Com relação ao \emph{tacto} diz S. Thomaz: « Cum sensus tactus sit quasi fundamentum aliorum sensuum, in organo uniuscujusque sensus non solum invenitur proprietas illius sensus, cujus est organum proprium, sed etiam proprietas tactus » (\emph{De} Ver., q. 22, a. 5).}. \textbar{}
